\chapter{A Concise Introduction to \LaTeX}\label{chap_LaTeX}\index{software!\LaTeX}

\section{Getting Started}
This template is intended to facilitate the adoption of \LaTeX{} for future homework submissions in UP504. This \texttt{.tex} file has been set up such that you can begin producing fairly advanced documents without getting too bogged down with coding (see appendix~\ref{sec_resources} for additional \LaTeX{} resources).

\subsection{Some Motivation}\label{sec_motivation}
The many advantages of using a typesetting system like \LaTeX{} arise from the fact that the computer takes care of all formatting issues, freeing you up to concentrate on ``the highest and best use'' of your time: delivering meaningful content. Unlike a conventional WYSIWYG\footnote{In WYSIWYG (What you see is what you get) text editing systems, the input and final output file are identical and are simultaneously displayed on screen. As you can see, footnotes are really easy too!} word processor such as MS Word, \LaTeX{} does not distract you by a clever reverse delegation of sectioning, pagination, object placement, numbering and other formatting tasks. Imagine being able to convert a letter-sized, one-column document with footnotes, tables and figures into an A4-sized, two-column layout with double-spaced, numbered lines and endnotes in less than 30 seconds! With \LaTeX{}, all you need is three easy commands and a few parameter changes to achieve that. The commands are \verb+\let\footnote=\endnote+, \verb+\doublespace+, and \verb+\linenumbers+ plus adding the terms \verb+twocolumn+ and \verb+a4+ to the \verb+\documentclass{}+ command.

When it comes to producing a professionally typeset document, there are of course many additional advantages to using \LaTeX{}, such as kerning (the proportional spacing between characters), the use of common ligatures (joining of one or more letters into a single glyph) or advanced hyphenation. You can find a comprehensive illustration of these aspects \href{http://nitens.org/taraborelli/latex}{here}.

\subsection{Some Basics}\label{sec_basics}
The input for \LaTeX{} is a plain text file that contains the text of the document, as well as the commands that tell \LaTeX{}
how to typeset the text. ``Whitespace'' characters, such as blanks or tabs, are treated uniformly as ``space'', several consecutive whitespace characters are treated as one space. \LaTeX{} treats a single line break as whitespace and an empty line between two lines of text defines the end of a paragraph. Several empty lines are treated the same as one empty line.

\subsubsection{Command Syntax}
\LaTeX{} commands are case sensitive, and take one of the following two formats:

\begin{itemize}
\item They start with a backslash \verb+\+ and then have a name consisting of letters only. Command names are terminated by a space, a number or any other `non-letter' character.
\item Alternatively, commands consist of a backslash and exactly one non-letter character.
\end{itemize}

In many instances, both command formats are available to achieve the same thing. For example, \verb+\newline+ or \verb+\\+ starts a new line without starting a new paragraph. \verb+\\*+ additionally prohibits a page break after the forced line break. Some commands require a parameter which has to be given between curly brackets \verb+{...}+ after the command name. Some commands take optional parameters which are inserted after the command name in square brackets \verb+[...]+. Generally, this implies the following command syntax:\\

\begin{center}
\verb+\+\texttt{command}\verb+[+\emph{optional parameter}\verb+]{+\emph{parameter}\verb+}+.
\end{center}


\subsubsection{Document Structure}
In order for \LaTeX{} input files to be compiled successfully, they have to follow a certain structure. Thus every input file must start with the command \verb+\documentclass{...}+ which defines what sort of document you intend to write. After that, you can define
commands to influence the style of the whole document, or load packages that add new features to the \LaTeX{} system. To load such a package you use the command \verb+\usepackage{...}+.

When all the setup work is done, you start the body of the text with the command \verb+\begin{document}+. Now you enter the text mixed with some useful \LaTeX{} commands. The area between \verb+\documentclass+ and \verb+\begin{document}+ is called the \emph{preamble}. The end of every input file is indicated by the \verb+\end{document}+ command. Anything that follows this command will be ignored by the \LaTeX{} compiler. For example, the simplest possible \LaTeX{} document is generated by the following code:

\begin{verbatim}
\documentclass{article}
\begin{document}
Small is beautiful.
\end{document}
\end{verbatim}

More info on document structure can be found \href{http://en.wikibooks.org/wiki/LaTeX/Document_Structure}{here}.

\subsection{Some Principles}\label{sec_principles}
When it comes to coding, \LaTeX{} is no different from any other interpreted language. Successful coding has a high fixed cost, a steep learning curve, but yields large economies to scale. Most questions about code are made more answerable by the addition of a minimal working example (MWE) or a short, self-contained, correct example (SSCCE). Your debugging question may have an obvious answer, or it may not. When enlisting the help of others in debugging your code, think about it from the point of view of the person whose help you're seeking. Thus you should always provide a block of code that can be copied, pasted directly into a file, compiled, tweaked, corrected, and pasted back.

In this context, ``minimal'' means that the problem is isolated and there is nothing in the sample that distracts from the error you have or the effect you desire. ``Working'' doesn't mean the problem is solved, it means the code sample can be copied-and-pasted and directly compiled.\footnote{The site \href{http://www.sscce.org}{sscce.org} is dedicated to the topic. There are also several articles on the topic as it applies specifically to LaTeX: \href{http://www.tex.ac.uk/cgi-bin/texfaq2html?label=minxampl}{``How to make a `minimum example'"} from the UK \TeX{} FAQ, \href{http://www.minimalbeispiel.de/mini-en.html}{``What is a minimal working example?"} from the \texttt{de.comp.text.tex} newsgroup.}

\subsection{Some Math}\label{sec_math}
One of the core strengths of \LaTeX{} is the typesetting of mathematical expressions. While the syntax takes some getting used to, you will soon discover its large economies to scale. Formulae are placed between the dollar sign operator, such as e.g. \verb+$\sum x^{2}_{i}$+ produces $\sum x^{2}_{i}$. Similarly, the mean of a discrete random variable $X$ with realisations $x_{1}, x_{2}, \ldots, x_{N}$ can be thought of as the probability-weighted sum of its realisation, i.e.

$$ \mu_{X} = \sum^{N}_{i=1}p_{i}x_{i}.$$

We can also write this same formula as an in-line text, $\mu_{X} = \sum^{N}_{i=1}p_{i}x_{i}$, without having to fret over possible implications for required spacing and formatting changes.

More complex formulae are -- with a little effort -- equally straightforward to typeset. For example, the probability density function for the normal distribution of a continuous random variable is

$$f(x) = \frac{1}{\sigma\sqrt{2\pi}}e^{-\frac{1}{2}\left(\frac{x-\mu}{\sigma}\right)^{2}}$$

You will be amazed how easy this is, once you practice a little. Instead of using the \verb+$+ or \verb+$$+ operators to invoke the math environment, we can also use the \verb+eqnarray+ environment for larger expressions. For example, we can define the mean square error (MSE) of an estimator $\hat{\theta}$ as

\begin{eqnarray}
\mathbb{E}\left[(\hat{\theta}-\theta)^2\right] &=& \mathbb{E}\left[(\hat{\theta}- \mu_{\hat{\theta}} + \mu_{\hat{\theta}} -\theta)^2\right]\\
&=&\mathbb{E}\left[(\hat{\theta}- \mu_{\hat{\theta}})^{2} + 2(\hat{\theta} - \mu_{\hat{\theta}})(\mu_{\hat{\theta}} - \theta) + (\mu_{\hat{\theta}} -\theta)^2\right]\nonumber\\
&=&\mathbb{E}\left[(\hat{\theta}- \mu_{\hat{\theta}})^{2}\right] + 2(\mu_{\hat{\theta}} - \theta)\mathbb{E}\left[\hat{\theta} - \mu_{\hat{\theta}}\right] + \mathbb{E}\Big[(\mu_{\hat{\theta}} -\theta)^2\Big]\label{eqn_mse1}\\
&=&\underbrace{\mathbb{E}\left[(\hat{\theta}- \mu_{\hat{\theta}})^{2}\right]}_{\textrm{variance of estimator } \hat{\theta}} + \underbrace{\mathbb{E}\Big[(\mu_{\hat{\theta}} -\theta)^2\Big]}_{\textrm{bias}^{2}}\label{eqn_mse2}
\end{eqnarray}

Notice that moving from equation~\eqref{eqn_mse1} to equation~\eqref{eqn_mse2}, the middle term vanishes because $\mathbb{E}\big[\hat{\theta} - \mu_{\hat{\theta}}\big]=\mathbb{E}\big[\hat{\theta}\big]-\mu_{\hat{\theta}}=0$. Thus, when choosing between two estimators, the MSE corresponds to a loss function that trades off the \emph{efficiency} of an estimator (its variance) against its \emph{unbiasedness}.

As you are stepping through the code that produces the above set of equations, you will notice another important aspect of \LaTeX{}; its internal reference and label handling. Objects such as equations, tables, figures or entire section headings can be referenced with ease by using the \verb+\label{+\emph{label\_name}\verb+}+ command first and then inserting the \verb+\ref{+\emph{label\_name}\verb+}+ or \verb+\eqref{+\emph{label\_name}\verb+}+ commands. This means that you will not have to worry about updating these references if your document changes or if you decide to apply a different numbering style (e.g. using roman numerals for section headings). The system will keep track of all labels and apply the appropriate numbering that is consistent with your formatting instructions. For example, this section is preceded by section~\ref{sec_principles} and followed by section~\ref{sec_tables}.


\subsection{Some Tables}\label{sec_tables}
The following provides a set of summary statistics for the data. This is presented in tables~\ref{tbl_Sum} and \ref{tbl_SumPlus}.

%%%%%%%%%%%%%%%%%%%%%%%%%%%%%%%%%%%%%%%%%%%%%%%%%%%%%%%%%%%%%%%%%%%%%
% Using a library such as "stargazer" or "xtable", you can
% simply use R to produce the LaTeX code for tables and then copy
% and paste that code here (NB: When overwriting the existing code
% you need to re-define the "\label{}" commands as they are not
% automatically produced by the R output).
%%%%%%%%%%%%%%%%%% TABLE CODE BEGINS %%%%%%%%%%%%%%%%%%%%%%%%%%%%%%%%
\begin{table}[htb] \centering
  \caption{Summary Statistics}
\footnotesize

\begin{tabular}{@{\extracolsep{5pt}}l c c c c c }\label{tbl_Sum}
\\[-1.8ex]\hline
\hline \\[-1.8ex]
Statistic & \multicolumn{1}{c}{\textit{N}} & \multicolumn{1}{c}{\textit{Mean}} & \multicolumn{1}{c}{\textit{St. Dev.}} & \multicolumn{1}{c}{\textit{Min}} & \multicolumn{1}{c}{\textit{Max}} \\
\hline \\[-1.8ex]
res\_vac & 1,554 & 103.627 & 113.150 & 0 & 871 \\
bus\_vac & 1,554 & 18.823 & 27.070 & 0 & 332 \\
oth\_vac & 1,554 & 0.045 & 0.297 & 0 & 7 \\
res\_tot & 1,554 & 1,525.257 & 732.713 & 0 & 7,563 \\
\hline \\[-1.8ex]
\normalsize
\end{tabular}
\end{table}
%%%%%%%%%%%%%%%%%% TABLE CODE ENDS %%%%%%%%%%%%%%%%%%%%%%%%%%%%%%%%


%%%%%%%%%%%%%%%%%% TABLE CODE BEGINS %%%%%%%%%%%%%%%%%%%%%%%%%%%%%%%%
\begin{table}[htp] \centering
  \caption{Summary Table with Additional Statistics}
\footnotesize

\begin{tabular}{@{\extracolsep{5pt}}l c c c c c c c c }\label{tbl_SumPlus}
\\[-1.8ex]\hline
\hline \\[-1.8ex]
Statistic & \multicolumn{1}{c}{\textit{N}} & \multicolumn{1}{c}{\textit{Mean}} & \multicolumn{1}{c}{\textit{St. Dev.}} & \multicolumn{1}{c}{\textit{Min}} & \multicolumn{1}{c}{\textit{Pctl(25)}} & \multicolumn{1}{c}{\textit{Median}} & \multicolumn{1}{c}{\textit{Pctl(75)}} & \multicolumn{1}{c}{\textit{Max}} \\
\hline \\[-1.8ex]
pop\_w & 1,554 & 2,501.398 & 1,787.798 & 0 & 1,123 & 2,473.5 & 3,569 & 11,632 \\
pop\_black & 1,554 & 732.153 & 1,109.979 & 0 & 33 & 173.5 & 953.5 & 5,687 \\
pop\_latino & 1,554 & 121.521 & 348.133 & 0 & 14 & 54 & 118 & 5,273 \\
hh\_median & 1,542 & 55,331.710 & 27,352.120 & 6,914 & 34,679 & 52,486.5 & 68,516.5 & 171,683 \\
\hline \\[-1.8ex]
\normalsize
\end{tabular}
\end{table}
%%%%%%%%%%%%%%%%%% TABLE CODE ENDS %%%%%%%%%%%%%%%%%%%%%%%%%%%%%%%%

\subsection{Some Figures}
Adding figures also does not take more than a couple of lines of code, using the \verb+\includegraphics{}+ command in the \verb+table+ environment. Best of all, once you have set this up in one document, you can simply adapt it for other documents. For example, figure~\ref{fig_Diversification} illustrates a positive correlation between city density and an index of relative industry diversity, calculated on the basis of all private employment at the two-digit NAICS level in 2010.

As you can see, the density-diversity relationship somewhat attenuates at the upper end of the city size distribution. This is partly because the population density in most larger cities is not exclusively driven by productivity differentials, but also by non-tradable amenity differences that both influence the quality-of-life and the quality of the business environment.

%%%%%%%%%%%%%%%%%% FIGURE CODE BEGINS %%%%%%%%%%%%%%%%%%%%%%%%%%%%%%
\begin{figure}[htp]
\caption{Relative industry diversity and population density of U.S. metropolitan areas, 2010}\label{fig_Diversification}
\begin{center}
\includegraphics[width=0.65\textwidth]{"Figures/RDI_PopDens_2010"}
\end{center}
\footnotesize
\parbox[b]{\textwidth}{\emph{Note}: The relative-diversity index (RDI) is computed by summing for each city $i$, over all $j$ sectors, the absolute value of the difference between each sector's share $s_{ij}$ in local employment and its share in national employment $s_{j}$. Formally, this yields $RDI_{i} = 1/ |\sum_{j}s_{ij} - s_{j}|$. Population density in thousands of inhabitants per square mile is plotted on a logarithmic scale.} \normalsize
\end{figure}
%%%%%%%%%%%%%%%%%% FIGURE CODE ENDS %%%%%%%%%%%%%%%%%%%%%%%%%%%%%%%%


\subsubsection{A Word on Placement}
Another nice thing about typesetting a document in this manner is that we do not have to worry about the specific placements of tables and figures. Notice that the expressions \verb+\begin{table}+ or \verb+\begin{figure}+ are followed by \verb+[htp]+, the so-called ``placement specifier'', which gives a preference for where the object, also called a ``float'', should be placed.\footnote{\texttt{[h]} means place the float where occurs in the text. \texttt{[t]} is the position at the top of the page; \texttt{[b]} denotes the bottom of the page, and \texttt{[p]} puts on a special page for floats only.} Indeed, depending on the amount of tables, figures and the general formatting constraints of your template, \LaTeX{} will figure out the optimal pagination.


\subsubsection{Some Additional Deep Insight}
Here we could have some additional text.


\section{Managing References}
In the \LaTeX{} environment, \href{http://www.bibtex.org/}{\textsc{Bib}\TeX{}} is the specific style for formatting lists of references and a \textsc{Bib}\TeX{} database is stored as a \texttt{.bib} file. It is a plain text file with a simple structure which can be viewed and edited easily. The file ``\texttt{References.bib}'' which you will find in the project folder ``\texttt{files}'' is an example of a minimal \textsc{Bib}\TeX{}.

\textsc{Bib}\TeX{} defines different entry types, with reports (\verb+@TECHREPORT{}+), articles (\verb+@ARTICLE{}+), books (\verb+@BOOK{}+) and book chapters (\verb+@INBOOK{}+) being the most common. While \textsc{Bib}\TeX{} files are simple enough to maintain using a text editor, it is more convenient to manage large \texttt{.bib} files with many entries via a reference management software that has a graphic user interface. Users of MS Word might be most familiar with the EndNotes software from Thompson Reuters for which a student version only costs \$104.99. I prefer to use software that is completely free, such as \href{http://www.mendeley.com/}{Mendeley} or, better still, \href{http://jabref.sourceforge.net/}{JabRef}, the system of choice in the open source community.


\subsection{In-text Citations}
In-text citations are really easy; we can refer to a specific piece of literature by simply inserting the respective labels for an entry that we have in our reference management database. The command \verb+\citet{+\emph{authorkey}\verb+}+ generates the most simple in-text citation. A variant of this command, \verb+\citep{}+, puts everything into parenthesis.

Notice that the specific look of the in-text citation depends on the citation style which is controlled with the command \verb+\bibliographystyle{+\emph{citationstyle}\verb+}+. The current documents uses the citation style \emph{econometrica}, the official citation style of the Econometric Society's flagship journal.\footnote{A comprehensive list of bibliography styles with examples can be found \href{http://www.mackichan.com/index.html?techtalk/632.htm~mainFrame}{here}.} The list of references is then automatically generated by the command \verb+\bibliography{+\emph{bibfile}\verb+}+ where the argument \emph{bibfile} is the name of your \texttt{.bib} file.

So, for example, I hope that you find the textbook by \citet{Dougherty2011} helpful. The article by \citet{Briant2010} might contain a surprise or two. \citet{Carrigan2012} and \citet{Goldstein2013} are useful reviews of the literature on the origins of the financial crisis.


\section{Conclusion}
No document is complete without it. Notice that there is a page break after the bibliography.

\newpage
\begin{subappendices}
\section{Additional \LaTeX{} Resources}\label{sec_resources}
There are a number of beginners' guides that I have posted on CTools. See \citet{Oetiker2011} for what is perhaps the most comprehensive introductory tutorial. Furthermore, a simple internet search with the key words ``\texttt{introduction}'' and ``\texttt{latex}'' will yield a plenty of accessible resources. Personally, I find \href{http://www.youtube.com/playlist?list=PLF975D9D3C9B50FF7}{this series} of Youtube videos particularly intuitive.

\section{My R Code}\label{sec_Rcode}
All of your homework submissions should contain the relevant \textsf{R} code that produces the key results in your work, such as tables, figures or estimation output. The \texttt{listings} package allows you to add non-formatted text as you would do with \verb+\begin{verbatim}+ but its main aim is to include the source code of any programming language within your document. The \texttt{listings} package supports highlighting of all the most common languages and it is highly customizable (see the coding set up in the preamble of the input file for this document).
\bigskip
% Copy and paste a cleaned-up version of your R code (plus relevant output) here, i.e. inside of the
% verbatim statements.
\begin{lstlisting}
setwd("[your directory structure]/UP_504/Data")

census <- read.csv("DA2_Census_TractsLarge.csv", header=T, na.strings="#N/A")
sub.census <- subset(census, census$metro>0)

library(stargazer)
library(xtable)

stargazer(sub.census[,2:7], title="Summary statistics")
stargazer(sub.census[c("pop_w", "pop_black", "pop_latino", "hh_median")],
    title="Summary table with additional statistics", median=T, iqr=T)

## Demonstration of longtable support
## Remember to insert \usepackage{longtable} on your LaTeX preamble

tbl.big <- xtable(sub.census[1:80,2:7],label="tbl.big",caption="Example of
    a Long Table Spanning Several Pages")
print(tbl.big,tabular.environment="longtable",include.colnames=T, floating=F)

## Demonstration of sidewaystable support
## Remember to insert \usepackage{rotating} on your LaTeX preamble

tbl.wide <- xtable(sub.census[1:15,2:12],label="tbl.wide",caption="A
    Wide Table")
print(tbl.wide,floating.environment="sidewaystable")
\end{lstlisting}
\normalsize

\section{A Very Long Table}
In some instances, you might need to display a large amount of data in one table that can span multiple pages or you might have a table that needs to be displayed in landscape format. The \texttt{longtable} environment is ideal for achieving the former and the \texttt{sidewaystable} environment deals with the latter task. Both functionalities are integrated nicely into the \texttt{xtable} library in \textsf{R}.

All the \textsf{R} code that was used to produce the tables in this document is printed in section~\ref{sec_Rcode} of this appendix.


%%%%%%%%%%%%%%%%%% TABLE CODE BEGINS %%%%%%%%%%%%%%%%%%%%%%%%%%%%%%%%
\begin{longtable}{rrrrrrr}
\caption{Example of a Long Table Spanning Several Pages}\label{tbl.big}\\
\hline
\hline
 & res\_vac & bus\_vac & oth\_vac & res\_tot & bus\_tot & oth\_tot \\
  \hline
\endfirsthead
 & res\_vac & bus\_vac & oth\_vac & res\_tot & bus\_tot & oth\_tot \\
  \hline
\endhead
28799 & 136 &   3 &   0 & 1307 &  36 &  27 \\
  28800 & 278 &   9 &   0 & 1216 &  44 &   2 \\
  28801 & 329 &  12 &   0 & 1480 &  38 &  14 \\
  28802 & 254 &  14 &   0 & 1144 &  56 &   2 \\
  28803 & 160 &  15 &   0 & 926 &  64 &   5 \\
  28804 & 353 &  14 &   0 & 1487 &  42 &   8 \\
  28805 & 302 &  19 &   0 & 1807 &  83 &  11 \\
  28806 & 267 &  14 &   0 & 824 &  43 &  15 \\
  28807 & 583 &  17 &   0 & 2103 &  77 &   4 \\
  28808 & 462 &  15 &   0 & 1504 &  41 &  24 \\
  28809 & 345 &  10 &   0 & 1517 &  72 &  70 \\
  28810 & 249 &   6 &   0 & 1518 &  23 &   2 \\
  28811 & 211 &  26 &   0 & 1512 & 102 &  14 \\
  28812 & 207 &  23 &   0 & 792 &  53 &  37 \\
  28813 & 526 &  11 &   0 & 1210 &  67 &  71 \\
  28814 & 273 &  17 &   0 & 2240 &  85 &  88 \\
  28815 & 218 &  14 &   0 & 1165 &  81 &  28 \\
  28816 & 259 &   7 &   0 & 1165 &  28 &   6 \\
  28817 & 145 &   7 &   0 & 1199 &  38 &   8 \\
  28818 & 219 &  18 &   0 & 896 &  88 &  12 \\
  28819 &   8 &  23 &   0 &  22 &  77 &   2 \\
  28820 & 489 &  37 &   0 & 1687 &  84 &  33 \\
  28821 & 255 &  26 &   0 & 1157 &  96 &  14 \\
  28822 & 167 &  22 &   0 & 1338 &  71 &  12 \\
  28823 &   7 &   1 &   0 & 698 &  23 &  51 \\
  28824 & 368 &  25 &   0 & 1609 &  89 &  21 \\
  28825 & 186 &  17 &   0 & 1762 &  86 &   6 \\
  28826 & 308 & 170 &   1 & 1129 & 854 & 242 \\
  28827 & 117 &  10 &   1 & 1297 &  49 &  98 \\
  28828 & 114 &  11 &   1 & 1648 &  57 &  39 \\
  28829 & 236 &  54 &   0 & 1238 & 266 & 101 \\
  28830 & 214 &  10 &   0 & 1268 &  77 &  15 \\
  28831 &  63 &  15 &   0 & 839 &  55 &  35 \\
  28832 & 373 &  21 &   0 & 1359 & 101 & 137 \\
  28833 & 159 &  20 &   0 & 1401 &  83 &  49 \\
  28834 & 277 &  23 &   0 & 2446 & 148 &  19 \\
  28835 & 267 &  26 &   0 & 1404 &  89 &  49 \\
  28836 & 157 &  25 &   0 & 771 & 122 &  30 \\
  28837 & 251 &  20 &   0 & 2511 & 113 &  46 \\
  28838 & 279 &  32 &   0 & 1745 & 102 &  14 \\
  28839 &   0 &   4 &   0 &  20 &  27 &   3 \\
  28840 &  68 &  25 &   0 & 1392 & 169 & 128 \\
  28841 &   6 &   0 &   0 & 565 &  41 &  22 \\
  28842 &  18 &   5 &   0 & 1117 &  55 &  22 \\
  28843 &  56 &  13 &   0 & 1158 & 105 &  43 \\
  28844 &  22 &   0 &   0 & 1255 &  48 &   8 \\
  28845 &  33 &   0 &   0 & 644 &  11 &   6 \\
  28846 &   1 &   0 &   0 & 767 &  29 &   7 \\
  28847 &   2 &   1 &   0 & 756 &  30 &   4 \\
  28848 & 336 &  34 &   0 & 1848 & 105 &  44 \\
  28849 & 231 &   8 &   0 & 2043 &  50 &  52 \\
  28850 &  78 &  16 &   0 & 1562 & 152 &  18 \\
  28851 &  30 &  24 &   0 & 1721 & 296 &  66 \\
  28852 &   0 &   0 &   0 & 986 &  30 &   5 \\
  28853 &  21 &   1 &   0 & 1917 &  75 &  11 \\
  28854 & 143 &  28 &   0 & 2611 & 187 &  89 \\
  28855 &  25 &   4 &   0 & 1202 &  43 &  14 \\
  28856 &  22 &   0 &   0 & 2302 &  27 &   5 \\
  28857 &  11 &   0 &   0 & 2214 &  17 &   0 \\
  28858 &  45 &   2 &   0 & 2111 &  23 &   6 \\
  28859 &   6 &   0 &   0 & 1153 &  65 &  32 \\
  28860 & 398 &  63 &   0 & 2103 & 229 &  33 \\
  28861 &  98 &  30 &   0 & 2261 & 318 & 175 \\
  28862 & 158 &  57 &   0 & 1852 & 404 & 137 \\
  28863 & 230 &  48 &   1 & 1918 & 244 & 189 \\
  28864 &  88 & 130 &   0 & 1528 & 842 & 135 \\
  28865 & 212 &  23 &   0 & 2681 & 151 &  20 \\
  28866 &  16 &  36 &   0 & 4280 & 277 &  83 \\
  28867 &   0 &   0 &   0 & 3154 & 142 & 201 \\
  28868 &  88 &  40 &   0 & 3177 & 269 & 221 \\
  28869 &  50 &  18 &   0 & 1096 & 168 &  98 \\
  28870 &  65 &   9 &   0 & 2697 & 148 & 261 \\
  28871 & 125 &  29 &   0 & 2672 & 221 &  36 \\
  28872 &  70 &   0 &   0 & 2578 &  25 &   4 \\
  28873 &   0 &   1 &   0 & 2465 &  22 &  80 \\
  28874 & 295 &  38 &   0 & 2903 & 194 &  34 \\
  28875 & 151 &  18 &   0 & 1566 & 153 &  32 \\
  28876 &  66 &  40 &   0 & 1495 & 233 &  24 \\
  28877 &  23 &   2 &   0 & 1813 &  16 &   0 \\
  28878 &  72 &  20 &   0 & 1496 & 135 &  32 \\
   \hline
\hline
\end{longtable}
%%%%%%%%%%%%%%%%%% TABLE CODE ENDS %%%%%%%%%%%%%%%%%%%%%%%%%%%%%%%%%%


%%%%%%%%%%%%%%%%%% TABLE CODE BEGINS %%%%%%%%%%%%%%%%%%%%%%%%%%%%%%%%
\begin{sidewaystable}[hbt]
\begin{center}
\begin{tabular}{rrrrrrrrrrrr}
  \hline
  \hline
 & res\_vac & bus\_vac & oth\_vac & res\_tot & bus\_tot & oth\_tot & avg\_vac\_r & avg\_vac\_b & avg\_vac\_o & metro & hh \\
  \hline
28799 & 136 &   3 &   0 & 1307 &  36 &  27 & 453.03 & 621.67 & 0.00 &   3 & 1218 \\
  28800 & 278 &   9 &   0 & 1216 &  44 &   2 & 758.30 & 823.33 & 0.00 &   3 & 898 \\
  28801 & 329 &  12 &   0 & 1480 &  38 &  14 & 493.99 & 668.83 & 0.00 &   3 & 1413 \\
  28802 & 254 &  14 &   0 & 1144 &  56 &   2 & 888.50 & 768.64 & 0.00 &   3 & 666 \\
  28803 & 160 &  15 &   0 & 926 &  64 &   5 & 769.26 & 1095.93 & 0.00 &   3 & 796 \\
  28804 & 353 &  14 &   0 & 1487 &  42 &   8 & 732.24 & 700.36 & 0.00 &   3 & 1251 \\
  28805 & 302 &  19 &   0 & 1807 &  83 &  11 & 439.34 & 810.47 & 0.00 &   3 & 1376 \\
  28806 & 267 &  14 &   0 & 824 &  43 &  15 & 782.27 & 1317.64 & 0.00 &   3 & 610 \\
  28807 & 583 &  17 &   0 & 2103 &  77 &   4 & 854.68 & 1177.06 & 0.00 &   3 & 1660 \\
  28808 & 462 &  15 &   0 & 1504 &  41 &  24 & 819.23 & 1138.00 & 0.00 &   3 & 1069 \\
  28809 & 345 &  10 &   0 & 1517 &  72 &  70 & 813.63 & 829.90 & 0.00 &   3 & 1078 \\
  28810 & 249 &   6 &   0 & 1518 &  23 &   2 & 539.71 & 1107.83 & 0.00 &   3 & 1346 \\
  28811 & 211 &  26 &   0 & 1512 & 102 &  14 & 483.38 & 1063.77 & 0.00 &   3 & 1425 \\
  28812 & 207 &  23 &   0 & 792 &  53 &  37 & 716.20 & 721.35 & 0.00 &   3 & 562 \\
  28813 & 526 &  11 &   0 & 1210 &  67 &  71 & 967.73 & 907.46 & 0.00 &   3 & 778 \\
   \hline
\end{tabular}
\caption{A Wide Table}
\label{tbl.wide}
\end{center}
\end{sidewaystable}
%%%%%%%%%%%%%%%%%% TABLE CODE ENDS %%%%%%%%%%%%%%%%%%%%%%%%%%%%%%%%%%

\section{Step-by-step installation instructions}
\begin{enumerate}
\item Visit \href{https://www.sharelatex.com}{\texttt{https://www.sharelatex.com}}, create an account and log in.
\item On CTools, go to the folder \texttt{UP 504 002 W13} $\rightarrow$ \texttt{Resources} $\rightarrow$ \texttt{LaTeX} $\rightarrow$ \texttt{Programs} and download the file \texttt{UP504\_LaTeXProject.zip} to a location on your machine that you can find again. Do not unzip the file.
\item In your \textsc{Share}\LaTeX{} environment, select the ``Projects'' item in the menu and click on the green \emph{New Project} button. Select ``Upload Zipped Project'' from the drop list and select the \texttt{.zip} folder that you have downloaded from CTools. Upload that folder.
\item Once the folder has been uploaded, you will see a screen with a section ``Collaborators'' and a section ``Options''. In the ``Root Document'' field in the ``Options'' section, select the \texttt{UP504\_HomeworkTemplate.tex}.
\item Click on the \texttt{UP504\_HomeworkTemplate.tex} icon in the vertical navigation pane on the left. You should now see the context of the input file. This is the file that gives \LaTeX{} all the necessary formatting instructions.
\item You can now compile this file into a \texttt{.pdf} output file by clicking on the green \emph{Recompile} button. Your compiled document is now visible. This file serves both as a template for future homeworks and as a manual that introduces the most essential functions for our class.
\item For future homeworks, you either simply replace the text in the template file or you can create a new blank project, copying and pasting the necessary code from the template.
\end{enumerate}

\renewcommand\bibname{References}
\bibliographystyle{econometrica}
\bibliography{Literature/OverLord}
\IfFileExists{Literature/IntroLaTeX.bbl}{\chapter{A Concise Introduction to \LaTeX}\label{chap_LaTeX}\index{software!\LaTeX}

\section{Getting Started}
This template is intended to facilitate the adoption of \LaTeX{} for future homework submissions in UP504. This \texttt{.tex} file has been set up such that you can begin producing fairly advanced documents without getting too bogged down with coding (see appendix~\ref{sec_resources} for additional \LaTeX{} resources).

\subsection{Some Motivation}\label{sec_motivation}
The many advantages of using a typesetting system like \LaTeX{} arise from the fact that the computer takes care of all formatting issues, freeing you up to concentrate on ``the highest and best use'' of your time: delivering meaningful content. Unlike a conventional WYSIWYG\footnote{In WYSIWYG (What you see is what you get) text editing systems, the input and final output file are identical and are simultaneously displayed on screen. As you can see, footnotes are really easy too!} word processor such as MS Word, \LaTeX{} does not distract you by a clever reverse delegation of sectioning, pagination, object placement, numbering and other formatting tasks. Imagine being able to convert a letter-sized, one-column document with footnotes, tables and figures into an A4-sized, two-column layout with double-spaced, numbered lines and endnotes in less than 30 seconds! With \LaTeX{}, all you need is three easy commands and a few parameter changes to achieve that. The commands are \verb+\let\footnote=\endnote+, \verb+\doublespace+, and \verb+\linenumbers+ plus adding the terms \verb+twocolumn+ and \verb+a4+ to the \verb+\documentclass{}+ command.

When it comes to producing a professionally typeset document, there are of course many additional advantages to using \LaTeX{}, such as kerning (the proportional spacing between characters), the use of common ligatures (joining of one or more letters into a single glyph) or advanced hyphenation. You can find a comprehensive illustration of these aspects \href{http://nitens.org/taraborelli/latex}{here}.

\subsection{Some Basics}\label{sec_basics}
The input for \LaTeX{} is a plain text file that contains the text of the document, as well as the commands that tell \LaTeX{}
how to typeset the text. ``Whitespace'' characters, such as blanks or tabs, are treated uniformly as ``space'', several consecutive whitespace characters are treated as one space. \LaTeX{} treats a single line break as whitespace and an empty line between two lines of text defines the end of a paragraph. Several empty lines are treated the same as one empty line.

\subsubsection{Command Syntax}
\LaTeX{} commands are case sensitive, and take one of the following two formats:

\begin{itemize}
\item They start with a backslash \verb+\+ and then have a name consisting of letters only. Command names are terminated by a space, a number or any other `non-letter' character.
\item Alternatively, commands consist of a backslash and exactly one non-letter character.
\end{itemize}

In many instances, both command formats are available to achieve the same thing. For example, \verb+\newline+ or \verb+\\+ starts a new line without starting a new paragraph. \verb+\\*+ additionally prohibits a page break after the forced line break. Some commands require a parameter which has to be given between curly brackets \verb+{...}+ after the command name. Some commands take optional parameters which are inserted after the command name in square brackets \verb+[...]+. Generally, this implies the following command syntax:\\

\begin{center}
\verb+\+\texttt{command}\verb+[+\emph{optional parameter}\verb+]{+\emph{parameter}\verb+}+.
\end{center}


\subsubsection{Document Structure}
In order for \LaTeX{} input files to be compiled successfully, they have to follow a certain structure. Thus every input file must start with the command \verb+\documentclass{...}+ which defines what sort of document you intend to write. After that, you can define
commands to influence the style of the whole document, or load packages that add new features to the \LaTeX{} system. To load such a package you use the command \verb+\usepackage{...}+.

When all the setup work is done, you start the body of the text with the command \verb+\begin{document}+. Now you enter the text mixed with some useful \LaTeX{} commands. The area between \verb+\documentclass+ and \verb+\begin{document}+ is called the \emph{preamble}. The end of every input file is indicated by the \verb+\end{document}+ command. Anything that follows this command will be ignored by the \LaTeX{} compiler. For example, the simplest possible \LaTeX{} document is generated by the following code:

\begin{verbatim}
\documentclass{article}
\begin{document}
Small is beautiful.
\end{document}
\end{verbatim}

More info on document structure can be found \href{http://en.wikibooks.org/wiki/LaTeX/Document_Structure}{here}.

\subsection{Some Principles}\label{sec_principles}
When it comes to coding, \LaTeX{} is no different from any other interpreted language. Successful coding has a high fixed cost, a steep learning curve, but yields large economies to scale. Most questions about code are made more answerable by the addition of a minimal working example (MWE) or a short, self-contained, correct example (SSCCE). Your debugging question may have an obvious answer, or it may not. When enlisting the help of others in debugging your code, think about it from the point of view of the person whose help you're seeking. Thus you should always provide a block of code that can be copied, pasted directly into a file, compiled, tweaked, corrected, and pasted back.

In this context, ``minimal'' means that the problem is isolated and there is nothing in the sample that distracts from the error you have or the effect you desire. ``Working'' doesn't mean the problem is solved, it means the code sample can be copied-and-pasted and directly compiled.\footnote{The site \href{http://www.sscce.org}{sscce.org} is dedicated to the topic. There are also several articles on the topic as it applies specifically to LaTeX: \href{http://www.tex.ac.uk/cgi-bin/texfaq2html?label=minxampl}{``How to make a `minimum example'"} from the UK \TeX{} FAQ, \href{http://www.minimalbeispiel.de/mini-en.html}{``What is a minimal working example?"} from the \texttt{de.comp.text.tex} newsgroup.}

\subsection{Some Math}\label{sec_math}
One of the core strengths of \LaTeX{} is the typesetting of mathematical expressions. While the syntax takes some getting used to, you will soon discover its large economies to scale. Formulae are placed between the dollar sign operator, such as e.g. \verb+$\sum x^{2}_{i}$+ produces $\sum x^{2}_{i}$. Similarly, the mean of a discrete random variable $X$ with realisations $x_{1}, x_{2}, \ldots, x_{N}$ can be thought of as the probability-weighted sum of its realisation, i.e.

$$ \mu_{X} = \sum^{N}_{i=1}p_{i}x_{i}.$$

We can also write this same formula as an in-line text, $\mu_{X} = \sum^{N}_{i=1}p_{i}x_{i}$, without having to fret over possible implications for required spacing and formatting changes.

More complex formulae are -- with a little effort -- equally straightforward to typeset. For example, the probability density function for the normal distribution of a continuous random variable is

$$f(x) = \frac{1}{\sigma\sqrt{2\pi}}e^{-\frac{1}{2}\left(\frac{x-\mu}{\sigma}\right)^{2}}$$

You will be amazed how easy this is, once you practice a little. Instead of using the \verb+$+ or \verb+$$+ operators to invoke the math environment, we can also use the \verb+eqnarray+ environment for larger expressions. For example, we can define the mean square error (MSE) of an estimator $\hat{\theta}$ as

\begin{eqnarray}
\mathbb{E}\left[(\hat{\theta}-\theta)^2\right] &=& \mathbb{E}\left[(\hat{\theta}- \mu_{\hat{\theta}} + \mu_{\hat{\theta}} -\theta)^2\right]\\
&=&\mathbb{E}\left[(\hat{\theta}- \mu_{\hat{\theta}})^{2} + 2(\hat{\theta} - \mu_{\hat{\theta}})(\mu_{\hat{\theta}} - \theta) + (\mu_{\hat{\theta}} -\theta)^2\right]\nonumber\\
&=&\mathbb{E}\left[(\hat{\theta}- \mu_{\hat{\theta}})^{2}\right] + 2(\mu_{\hat{\theta}} - \theta)\mathbb{E}\left[\hat{\theta} - \mu_{\hat{\theta}}\right] + \mathbb{E}\Big[(\mu_{\hat{\theta}} -\theta)^2\Big]\label{eqn_mse1}\\
&=&\underbrace{\mathbb{E}\left[(\hat{\theta}- \mu_{\hat{\theta}})^{2}\right]}_{\textrm{variance of estimator } \hat{\theta}} + \underbrace{\mathbb{E}\Big[(\mu_{\hat{\theta}} -\theta)^2\Big]}_{\textrm{bias}^{2}}\label{eqn_mse2}
\end{eqnarray}

Notice that moving from equation~\eqref{eqn_mse1} to equation~\eqref{eqn_mse2}, the middle term vanishes because $\mathbb{E}\big[\hat{\theta} - \mu_{\hat{\theta}}\big]=\mathbb{E}\big[\hat{\theta}\big]-\mu_{\hat{\theta}}=0$. Thus, when choosing between two estimators, the MSE corresponds to a loss function that trades off the \emph{efficiency} of an estimator (its variance) against its \emph{unbiasedness}.

As you are stepping through the code that produces the above set of equations, you will notice another important aspect of \LaTeX{}; its internal reference and label handling. Objects such as equations, tables, figures or entire section headings can be referenced with ease by using the \verb+\label{+\emph{label\_name}\verb+}+ command first and then inserting the \verb+\ref{+\emph{label\_name}\verb+}+ or \verb+\eqref{+\emph{label\_name}\verb+}+ commands. This means that you will not have to worry about updating these references if your document changes or if you decide to apply a different numbering style (e.g. using roman numerals for section headings). The system will keep track of all labels and apply the appropriate numbering that is consistent with your formatting instructions. For example, this section is preceded by section~\ref{sec_principles} and followed by section~\ref{sec_tables}.


\subsection{Some Tables}\label{sec_tables}
The following provides a set of summary statistics for the data. This is presented in tables~\ref{tbl_Sum} and \ref{tbl_SumPlus}.

%%%%%%%%%%%%%%%%%%%%%%%%%%%%%%%%%%%%%%%%%%%%%%%%%%%%%%%%%%%%%%%%%%%%%
% Using a library such as "stargazer" or "xtable", you can
% simply use R to produce the LaTeX code for tables and then copy
% and paste that code here (NB: When overwriting the existing code
% you need to re-define the "\label{}" commands as they are not
% automatically produced by the R output).
%%%%%%%%%%%%%%%%%% TABLE CODE BEGINS %%%%%%%%%%%%%%%%%%%%%%%%%%%%%%%%
\begin{table}[htb] \centering
  \caption{Summary Statistics}
\footnotesize

\begin{tabular}{@{\extracolsep{5pt}}l c c c c c }\label{tbl_Sum}
\\[-1.8ex]\hline
\hline \\[-1.8ex]
Statistic & \multicolumn{1}{c}{\textit{N}} & \multicolumn{1}{c}{\textit{Mean}} & \multicolumn{1}{c}{\textit{St. Dev.}} & \multicolumn{1}{c}{\textit{Min}} & \multicolumn{1}{c}{\textit{Max}} \\
\hline \\[-1.8ex]
res\_vac & 1,554 & 103.627 & 113.150 & 0 & 871 \\
bus\_vac & 1,554 & 18.823 & 27.070 & 0 & 332 \\
oth\_vac & 1,554 & 0.045 & 0.297 & 0 & 7 \\
res\_tot & 1,554 & 1,525.257 & 732.713 & 0 & 7,563 \\
\hline \\[-1.8ex]
\normalsize
\end{tabular}
\end{table}
%%%%%%%%%%%%%%%%%% TABLE CODE ENDS %%%%%%%%%%%%%%%%%%%%%%%%%%%%%%%%


%%%%%%%%%%%%%%%%%% TABLE CODE BEGINS %%%%%%%%%%%%%%%%%%%%%%%%%%%%%%%%
\begin{table}[htp] \centering
  \caption{Summary Table with Additional Statistics}
\footnotesize

\begin{tabular}{@{\extracolsep{5pt}}l c c c c c c c c }\label{tbl_SumPlus}
\\[-1.8ex]\hline
\hline \\[-1.8ex]
Statistic & \multicolumn{1}{c}{\textit{N}} & \multicolumn{1}{c}{\textit{Mean}} & \multicolumn{1}{c}{\textit{St. Dev.}} & \multicolumn{1}{c}{\textit{Min}} & \multicolumn{1}{c}{\textit{Pctl(25)}} & \multicolumn{1}{c}{\textit{Median}} & \multicolumn{1}{c}{\textit{Pctl(75)}} & \multicolumn{1}{c}{\textit{Max}} \\
\hline \\[-1.8ex]
pop\_w & 1,554 & 2,501.398 & 1,787.798 & 0 & 1,123 & 2,473.5 & 3,569 & 11,632 \\
pop\_black & 1,554 & 732.153 & 1,109.979 & 0 & 33 & 173.5 & 953.5 & 5,687 \\
pop\_latino & 1,554 & 121.521 & 348.133 & 0 & 14 & 54 & 118 & 5,273 \\
hh\_median & 1,542 & 55,331.710 & 27,352.120 & 6,914 & 34,679 & 52,486.5 & 68,516.5 & 171,683 \\
\hline \\[-1.8ex]
\normalsize
\end{tabular}
\end{table}
%%%%%%%%%%%%%%%%%% TABLE CODE ENDS %%%%%%%%%%%%%%%%%%%%%%%%%%%%%%%%

\subsection{Some Figures}
Adding figures also does not take more than a couple of lines of code, using the \verb+\includegraphics{}+ command in the \verb+table+ environment. Best of all, once you have set this up in one document, you can simply adapt it for other documents. For example, figure~\ref{fig_Diversification} illustrates a positive correlation between city density and an index of relative industry diversity, calculated on the basis of all private employment at the two-digit NAICS level in 2010.

As you can see, the density-diversity relationship somewhat attenuates at the upper end of the city size distribution. This is partly because the population density in most larger cities is not exclusively driven by productivity differentials, but also by non-tradable amenity differences that both influence the quality-of-life and the quality of the business environment.

%%%%%%%%%%%%%%%%%% FIGURE CODE BEGINS %%%%%%%%%%%%%%%%%%%%%%%%%%%%%%
\begin{figure}[htp]
\caption{Relative industry diversity and population density of U.S. metropolitan areas, 2010}\label{fig_Diversification}
\begin{center}
\includegraphics[width=0.65\textwidth]{"Figures/RDI_PopDens_2010"}
\end{center}
\footnotesize
\parbox[b]{\textwidth}{\emph{Note}: The relative-diversity index (RDI) is computed by summing for each city $i$, over all $j$ sectors, the absolute value of the difference between each sector's share $s_{ij}$ in local employment and its share in national employment $s_{j}$. Formally, this yields $RDI_{i} = 1/ |\sum_{j}s_{ij} - s_{j}|$. Population density in thousands of inhabitants per square mile is plotted on a logarithmic scale.} \normalsize
\end{figure}
%%%%%%%%%%%%%%%%%% FIGURE CODE ENDS %%%%%%%%%%%%%%%%%%%%%%%%%%%%%%%%


\subsubsection{A Word on Placement}
Another nice thing about typesetting a document in this manner is that we do not have to worry about the specific placements of tables and figures. Notice that the expressions \verb+\begin{table}+ or \verb+\begin{figure}+ are followed by \verb+[htp]+, the so-called ``placement specifier'', which gives a preference for where the object, also called a ``float'', should be placed.\footnote{\texttt{[h]} means place the float where occurs in the text. \texttt{[t]} is the position at the top of the page; \texttt{[b]} denotes the bottom of the page, and \texttt{[p]} puts on a special page for floats only.} Indeed, depending on the amount of tables, figures and the general formatting constraints of your template, \LaTeX{} will figure out the optimal pagination.


\subsubsection{Some Additional Deep Insight}
Here we could have some additional text.


\section{Managing References}
In the \LaTeX{} environment, \href{http://www.bibtex.org/}{\textsc{Bib}\TeX{}} is the specific style for formatting lists of references and a \textsc{Bib}\TeX{} database is stored as a \texttt{.bib} file. It is a plain text file with a simple structure which can be viewed and edited easily. The file ``\texttt{References.bib}'' which you will find in the project folder ``\texttt{files}'' is an example of a minimal \textsc{Bib}\TeX{}.

\textsc{Bib}\TeX{} defines different entry types, with reports (\verb+@TECHREPORT{}+), articles (\verb+@ARTICLE{}+), books (\verb+@BOOK{}+) and book chapters (\verb+@INBOOK{}+) being the most common. While \textsc{Bib}\TeX{} files are simple enough to maintain using a text editor, it is more convenient to manage large \texttt{.bib} files with many entries via a reference management software that has a graphic user interface. Users of MS Word might be most familiar with the EndNotes software from Thompson Reuters for which a student version only costs \$104.99. I prefer to use software that is completely free, such as \href{http://www.mendeley.com/}{Mendeley} or, better still, \href{http://jabref.sourceforge.net/}{JabRef}, the system of choice in the open source community.


\subsection{In-text Citations}
In-text citations are really easy; we can refer to a specific piece of literature by simply inserting the respective labels for an entry that we have in our reference management database. The command \verb+\citet{+\emph{authorkey}\verb+}+ generates the most simple in-text citation. A variant of this command, \verb+\citep{}+, puts everything into parenthesis.

Notice that the specific look of the in-text citation depends on the citation style which is controlled with the command \verb+\bibliographystyle{+\emph{citationstyle}\verb+}+. The current documents uses the citation style \emph{econometrica}, the official citation style of the Econometric Society's flagship journal.\footnote{A comprehensive list of bibliography styles with examples can be found \href{http://www.mackichan.com/index.html?techtalk/632.htm~mainFrame}{here}.} The list of references is then automatically generated by the command \verb+\bibliography{+\emph{bibfile}\verb+}+ where the argument \emph{bibfile} is the name of your \texttt{.bib} file.

So, for example, I hope that you find the textbook by \citet{Dougherty2011} helpful. The article by \citet{Briant2010} might contain a surprise or two. \citet{Carrigan2012} and \citet{Goldstein2013} are useful reviews of the literature on the origins of the financial crisis.


\section{Conclusion}
No document is complete without it. Notice that there is a page break after the bibliography.

\newpage
\begin{subappendices}
\section{Additional \LaTeX{} Resources}\label{sec_resources}
There are a number of beginners' guides that I have posted on CTools. See \citet{Oetiker2011} for what is perhaps the most comprehensive introductory tutorial. Furthermore, a simple internet search with the key words ``\texttt{introduction}'' and ``\texttt{latex}'' will yield a plenty of accessible resources. Personally, I find \href{http://www.youtube.com/playlist?list=PLF975D9D3C9B50FF7}{this series} of Youtube videos particularly intuitive.

\section{My R Code}\label{sec_Rcode}
All of your homework submissions should contain the relevant \textsf{R} code that produces the key results in your work, such as tables, figures or estimation output. The \texttt{listings} package allows you to add non-formatted text as you would do with \verb+\begin{verbatim}+ but its main aim is to include the source code of any programming language within your document. The \texttt{listings} package supports highlighting of all the most common languages and it is highly customizable (see the coding set up in the preamble of the input file for this document).
\bigskip
% Copy and paste a cleaned-up version of your R code (plus relevant output) here, i.e. inside of the
% verbatim statements.
\begin{lstlisting}
setwd("[your directory structure]/UP_504/Data")

census <- read.csv("DA2_Census_TractsLarge.csv", header=T, na.strings="#N/A")
sub.census <- subset(census, census$metro>0)

library(stargazer)
library(xtable)

stargazer(sub.census[,2:7], title="Summary statistics")
stargazer(sub.census[c("pop_w", "pop_black", "pop_latino", "hh_median")],
    title="Summary table with additional statistics", median=T, iqr=T)

## Demonstration of longtable support
## Remember to insert \usepackage{longtable} on your LaTeX preamble

tbl.big <- xtable(sub.census[1:80,2:7],label="tbl.big",caption="Example of
    a Long Table Spanning Several Pages")
print(tbl.big,tabular.environment="longtable",include.colnames=T, floating=F)

## Demonstration of sidewaystable support
## Remember to insert \usepackage{rotating} on your LaTeX preamble

tbl.wide <- xtable(sub.census[1:15,2:12],label="tbl.wide",caption="A
    Wide Table")
print(tbl.wide,floating.environment="sidewaystable")
\end{lstlisting}
\normalsize

\section{A Very Long Table}
In some instances, you might need to display a large amount of data in one table that can span multiple pages or you might have a table that needs to be displayed in landscape format. The \texttt{longtable} environment is ideal for achieving the former and the \texttt{sidewaystable} environment deals with the latter task. Both functionalities are integrated nicely into the \texttt{xtable} library in \textsf{R}.

All the \textsf{R} code that was used to produce the tables in this document is printed in section~\ref{sec_Rcode} of this appendix.


%%%%%%%%%%%%%%%%%% TABLE CODE BEGINS %%%%%%%%%%%%%%%%%%%%%%%%%%%%%%%%
\begin{longtable}{rrrrrrr}
\caption{Example of a Long Table Spanning Several Pages}\label{tbl.big}\\
\hline
\hline
 & res\_vac & bus\_vac & oth\_vac & res\_tot & bus\_tot & oth\_tot \\
  \hline
\endfirsthead
 & res\_vac & bus\_vac & oth\_vac & res\_tot & bus\_tot & oth\_tot \\
  \hline
\endhead
28799 & 136 &   3 &   0 & 1307 &  36 &  27 \\
  28800 & 278 &   9 &   0 & 1216 &  44 &   2 \\
  28801 & 329 &  12 &   0 & 1480 &  38 &  14 \\
  28802 & 254 &  14 &   0 & 1144 &  56 &   2 \\
  28803 & 160 &  15 &   0 & 926 &  64 &   5 \\
  28804 & 353 &  14 &   0 & 1487 &  42 &   8 \\
  28805 & 302 &  19 &   0 & 1807 &  83 &  11 \\
  28806 & 267 &  14 &   0 & 824 &  43 &  15 \\
  28807 & 583 &  17 &   0 & 2103 &  77 &   4 \\
  28808 & 462 &  15 &   0 & 1504 &  41 &  24 \\
  28809 & 345 &  10 &   0 & 1517 &  72 &  70 \\
  28810 & 249 &   6 &   0 & 1518 &  23 &   2 \\
  28811 & 211 &  26 &   0 & 1512 & 102 &  14 \\
  28812 & 207 &  23 &   0 & 792 &  53 &  37 \\
  28813 & 526 &  11 &   0 & 1210 &  67 &  71 \\
  28814 & 273 &  17 &   0 & 2240 &  85 &  88 \\
  28815 & 218 &  14 &   0 & 1165 &  81 &  28 \\
  28816 & 259 &   7 &   0 & 1165 &  28 &   6 \\
  28817 & 145 &   7 &   0 & 1199 &  38 &   8 \\
  28818 & 219 &  18 &   0 & 896 &  88 &  12 \\
  28819 &   8 &  23 &   0 &  22 &  77 &   2 \\
  28820 & 489 &  37 &   0 & 1687 &  84 &  33 \\
  28821 & 255 &  26 &   0 & 1157 &  96 &  14 \\
  28822 & 167 &  22 &   0 & 1338 &  71 &  12 \\
  28823 &   7 &   1 &   0 & 698 &  23 &  51 \\
  28824 & 368 &  25 &   0 & 1609 &  89 &  21 \\
  28825 & 186 &  17 &   0 & 1762 &  86 &   6 \\
  28826 & 308 & 170 &   1 & 1129 & 854 & 242 \\
  28827 & 117 &  10 &   1 & 1297 &  49 &  98 \\
  28828 & 114 &  11 &   1 & 1648 &  57 &  39 \\
  28829 & 236 &  54 &   0 & 1238 & 266 & 101 \\
  28830 & 214 &  10 &   0 & 1268 &  77 &  15 \\
  28831 &  63 &  15 &   0 & 839 &  55 &  35 \\
  28832 & 373 &  21 &   0 & 1359 & 101 & 137 \\
  28833 & 159 &  20 &   0 & 1401 &  83 &  49 \\
  28834 & 277 &  23 &   0 & 2446 & 148 &  19 \\
  28835 & 267 &  26 &   0 & 1404 &  89 &  49 \\
  28836 & 157 &  25 &   0 & 771 & 122 &  30 \\
  28837 & 251 &  20 &   0 & 2511 & 113 &  46 \\
  28838 & 279 &  32 &   0 & 1745 & 102 &  14 \\
  28839 &   0 &   4 &   0 &  20 &  27 &   3 \\
  28840 &  68 &  25 &   0 & 1392 & 169 & 128 \\
  28841 &   6 &   0 &   0 & 565 &  41 &  22 \\
  28842 &  18 &   5 &   0 & 1117 &  55 &  22 \\
  28843 &  56 &  13 &   0 & 1158 & 105 &  43 \\
  28844 &  22 &   0 &   0 & 1255 &  48 &   8 \\
  28845 &  33 &   0 &   0 & 644 &  11 &   6 \\
  28846 &   1 &   0 &   0 & 767 &  29 &   7 \\
  28847 &   2 &   1 &   0 & 756 &  30 &   4 \\
  28848 & 336 &  34 &   0 & 1848 & 105 &  44 \\
  28849 & 231 &   8 &   0 & 2043 &  50 &  52 \\
  28850 &  78 &  16 &   0 & 1562 & 152 &  18 \\
  28851 &  30 &  24 &   0 & 1721 & 296 &  66 \\
  28852 &   0 &   0 &   0 & 986 &  30 &   5 \\
  28853 &  21 &   1 &   0 & 1917 &  75 &  11 \\
  28854 & 143 &  28 &   0 & 2611 & 187 &  89 \\
  28855 &  25 &   4 &   0 & 1202 &  43 &  14 \\
  28856 &  22 &   0 &   0 & 2302 &  27 &   5 \\
  28857 &  11 &   0 &   0 & 2214 &  17 &   0 \\
  28858 &  45 &   2 &   0 & 2111 &  23 &   6 \\
  28859 &   6 &   0 &   0 & 1153 &  65 &  32 \\
  28860 & 398 &  63 &   0 & 2103 & 229 &  33 \\
  28861 &  98 &  30 &   0 & 2261 & 318 & 175 \\
  28862 & 158 &  57 &   0 & 1852 & 404 & 137 \\
  28863 & 230 &  48 &   1 & 1918 & 244 & 189 \\
  28864 &  88 & 130 &   0 & 1528 & 842 & 135 \\
  28865 & 212 &  23 &   0 & 2681 & 151 &  20 \\
  28866 &  16 &  36 &   0 & 4280 & 277 &  83 \\
  28867 &   0 &   0 &   0 & 3154 & 142 & 201 \\
  28868 &  88 &  40 &   0 & 3177 & 269 & 221 \\
  28869 &  50 &  18 &   0 & 1096 & 168 &  98 \\
  28870 &  65 &   9 &   0 & 2697 & 148 & 261 \\
  28871 & 125 &  29 &   0 & 2672 & 221 &  36 \\
  28872 &  70 &   0 &   0 & 2578 &  25 &   4 \\
  28873 &   0 &   1 &   0 & 2465 &  22 &  80 \\
  28874 & 295 &  38 &   0 & 2903 & 194 &  34 \\
  28875 & 151 &  18 &   0 & 1566 & 153 &  32 \\
  28876 &  66 &  40 &   0 & 1495 & 233 &  24 \\
  28877 &  23 &   2 &   0 & 1813 &  16 &   0 \\
  28878 &  72 &  20 &   0 & 1496 & 135 &  32 \\
   \hline
\hline
\end{longtable}
%%%%%%%%%%%%%%%%%% TABLE CODE ENDS %%%%%%%%%%%%%%%%%%%%%%%%%%%%%%%%%%


%%%%%%%%%%%%%%%%%% TABLE CODE BEGINS %%%%%%%%%%%%%%%%%%%%%%%%%%%%%%%%
\begin{sidewaystable}[hbt]
\begin{center}
\begin{tabular}{rrrrrrrrrrrr}
  \hline
  \hline
 & res\_vac & bus\_vac & oth\_vac & res\_tot & bus\_tot & oth\_tot & avg\_vac\_r & avg\_vac\_b & avg\_vac\_o & metro & hh \\
  \hline
28799 & 136 &   3 &   0 & 1307 &  36 &  27 & 453.03 & 621.67 & 0.00 &   3 & 1218 \\
  28800 & 278 &   9 &   0 & 1216 &  44 &   2 & 758.30 & 823.33 & 0.00 &   3 & 898 \\
  28801 & 329 &  12 &   0 & 1480 &  38 &  14 & 493.99 & 668.83 & 0.00 &   3 & 1413 \\
  28802 & 254 &  14 &   0 & 1144 &  56 &   2 & 888.50 & 768.64 & 0.00 &   3 & 666 \\
  28803 & 160 &  15 &   0 & 926 &  64 &   5 & 769.26 & 1095.93 & 0.00 &   3 & 796 \\
  28804 & 353 &  14 &   0 & 1487 &  42 &   8 & 732.24 & 700.36 & 0.00 &   3 & 1251 \\
  28805 & 302 &  19 &   0 & 1807 &  83 &  11 & 439.34 & 810.47 & 0.00 &   3 & 1376 \\
  28806 & 267 &  14 &   0 & 824 &  43 &  15 & 782.27 & 1317.64 & 0.00 &   3 & 610 \\
  28807 & 583 &  17 &   0 & 2103 &  77 &   4 & 854.68 & 1177.06 & 0.00 &   3 & 1660 \\
  28808 & 462 &  15 &   0 & 1504 &  41 &  24 & 819.23 & 1138.00 & 0.00 &   3 & 1069 \\
  28809 & 345 &  10 &   0 & 1517 &  72 &  70 & 813.63 & 829.90 & 0.00 &   3 & 1078 \\
  28810 & 249 &   6 &   0 & 1518 &  23 &   2 & 539.71 & 1107.83 & 0.00 &   3 & 1346 \\
  28811 & 211 &  26 &   0 & 1512 & 102 &  14 & 483.38 & 1063.77 & 0.00 &   3 & 1425 \\
  28812 & 207 &  23 &   0 & 792 &  53 &  37 & 716.20 & 721.35 & 0.00 &   3 & 562 \\
  28813 & 526 &  11 &   0 & 1210 &  67 &  71 & 967.73 & 907.46 & 0.00 &   3 & 778 \\
   \hline
\end{tabular}
\caption{A Wide Table}
\label{tbl.wide}
\end{center}
\end{sidewaystable}
%%%%%%%%%%%%%%%%%% TABLE CODE ENDS %%%%%%%%%%%%%%%%%%%%%%%%%%%%%%%%%%

\section{Step-by-step installation instructions}
\begin{enumerate}
\item Visit \href{https://www.sharelatex.com}{\texttt{https://www.sharelatex.com}}, create an account and log in.
\item On CTools, go to the folder \texttt{UP 504 002 W13} $\rightarrow$ \texttt{Resources} $\rightarrow$ \texttt{LaTeX} $\rightarrow$ \texttt{Programs} and download the file \texttt{UP504\_LaTeXProject.zip} to a location on your machine that you can find again. Do not unzip the file.
\item In your \textsc{Share}\LaTeX{} environment, select the ``Projects'' item in the menu and click on the green \emph{New Project} button. Select ``Upload Zipped Project'' from the drop list and select the \texttt{.zip} folder that you have downloaded from CTools. Upload that folder.
\item Once the folder has been uploaded, you will see a screen with a section ``Collaborators'' and a section ``Options''. In the ``Root Document'' field in the ``Options'' section, select the \texttt{UP504\_HomeworkTemplate.tex}.
\item Click on the \texttt{UP504\_HomeworkTemplate.tex} icon in the vertical navigation pane on the left. You should now see the context of the input file. This is the file that gives \LaTeX{} all the necessary formatting instructions.
\item You can now compile this file into a \texttt{.pdf} output file by clicking on the green \emph{Recompile} button. Your compiled document is now visible. This file serves both as a template for future homeworks and as a manual that introduces the most essential functions for our class.
\item For future homeworks, you either simply replace the text in the template file or you can create a new blank project, copying and pasting the necessary code from the template.
\end{enumerate}

\renewcommand\bibname{References}
\bibliographystyle{econometrica}
\bibliography{Literature/OverLord}
\IfFileExists{Literature/IntroLaTeX.bbl}{\chapter{A Concise Introduction to \LaTeX}\label{chap_LaTeX}\index{software!\LaTeX}

\section{Getting Started}
This template is intended to facilitate the adoption of \LaTeX{} for future homework submissions in UP504. This \texttt{.tex} file has been set up such that you can begin producing fairly advanced documents without getting too bogged down with coding (see appendix~\ref{sec_resources} for additional \LaTeX{} resources).

\subsection{Some Motivation}\label{sec_motivation}
The many advantages of using a typesetting system like \LaTeX{} arise from the fact that the computer takes care of all formatting issues, freeing you up to concentrate on ``the highest and best use'' of your time: delivering meaningful content. Unlike a conventional WYSIWYG\footnote{In WYSIWYG (What you see is what you get) text editing systems, the input and final output file are identical and are simultaneously displayed on screen. As you can see, footnotes are really easy too!} word processor such as MS Word, \LaTeX{} does not distract you by a clever reverse delegation of sectioning, pagination, object placement, numbering and other formatting tasks. Imagine being able to convert a letter-sized, one-column document with footnotes, tables and figures into an A4-sized, two-column layout with double-spaced, numbered lines and endnotes in less than 30 seconds! With \LaTeX{}, all you need is three easy commands and a few parameter changes to achieve that. The commands are \verb+\let\footnote=\endnote+, \verb+\doublespace+, and \verb+\linenumbers+ plus adding the terms \verb+twocolumn+ and \verb+a4+ to the \verb+\documentclass{}+ command.

When it comes to producing a professionally typeset document, there are of course many additional advantages to using \LaTeX{}, such as kerning (the proportional spacing between characters), the use of common ligatures (joining of one or more letters into a single glyph) or advanced hyphenation. You can find a comprehensive illustration of these aspects \href{http://nitens.org/taraborelli/latex}{here}.

\subsection{Some Basics}\label{sec_basics}
The input for \LaTeX{} is a plain text file that contains the text of the document, as well as the commands that tell \LaTeX{}
how to typeset the text. ``Whitespace'' characters, such as blanks or tabs, are treated uniformly as ``space'', several consecutive whitespace characters are treated as one space. \LaTeX{} treats a single line break as whitespace and an empty line between two lines of text defines the end of a paragraph. Several empty lines are treated the same as one empty line.

\subsubsection{Command Syntax}
\LaTeX{} commands are case sensitive, and take one of the following two formats:

\begin{itemize}
\item They start with a backslash \verb+\+ and then have a name consisting of letters only. Command names are terminated by a space, a number or any other `non-letter' character.
\item Alternatively, commands consist of a backslash and exactly one non-letter character.
\end{itemize}

In many instances, both command formats are available to achieve the same thing. For example, \verb+\newline+ or \verb+\\+ starts a new line without starting a new paragraph. \verb+\\*+ additionally prohibits a page break after the forced line break. Some commands require a parameter which has to be given between curly brackets \verb+{...}+ after the command name. Some commands take optional parameters which are inserted after the command name in square brackets \verb+[...]+. Generally, this implies the following command syntax:\\

\begin{center}
\verb+\+\texttt{command}\verb+[+\emph{optional parameter}\verb+]{+\emph{parameter}\verb+}+.
\end{center}


\subsubsection{Document Structure}
In order for \LaTeX{} input files to be compiled successfully, they have to follow a certain structure. Thus every input file must start with the command \verb+\documentclass{...}+ which defines what sort of document you intend to write. After that, you can define
commands to influence the style of the whole document, or load packages that add new features to the \LaTeX{} system. To load such a package you use the command \verb+\usepackage{...}+.

When all the setup work is done, you start the body of the text with the command \verb+\begin{document}+. Now you enter the text mixed with some useful \LaTeX{} commands. The area between \verb+\documentclass+ and \verb+\begin{document}+ is called the \emph{preamble}. The end of every input file is indicated by the \verb+\end{document}+ command. Anything that follows this command will be ignored by the \LaTeX{} compiler. For example, the simplest possible \LaTeX{} document is generated by the following code:

\begin{verbatim}
\documentclass{article}
\begin{document}
Small is beautiful.
\end{document}
\end{verbatim}

More info on document structure can be found \href{http://en.wikibooks.org/wiki/LaTeX/Document_Structure}{here}.

\subsection{Some Principles}\label{sec_principles}
When it comes to coding, \LaTeX{} is no different from any other interpreted language. Successful coding has a high fixed cost, a steep learning curve, but yields large economies to scale. Most questions about code are made more answerable by the addition of a minimal working example (MWE) or a short, self-contained, correct example (SSCCE). Your debugging question may have an obvious answer, or it may not. When enlisting the help of others in debugging your code, think about it from the point of view of the person whose help you're seeking. Thus you should always provide a block of code that can be copied, pasted directly into a file, compiled, tweaked, corrected, and pasted back.

In this context, ``minimal'' means that the problem is isolated and there is nothing in the sample that distracts from the error you have or the effect you desire. ``Working'' doesn't mean the problem is solved, it means the code sample can be copied-and-pasted and directly compiled.\footnote{The site \href{http://www.sscce.org}{sscce.org} is dedicated to the topic. There are also several articles on the topic as it applies specifically to LaTeX: \href{http://www.tex.ac.uk/cgi-bin/texfaq2html?label=minxampl}{``How to make a `minimum example'"} from the UK \TeX{} FAQ, \href{http://www.minimalbeispiel.de/mini-en.html}{``What is a minimal working example?"} from the \texttt{de.comp.text.tex} newsgroup.}

\subsection{Some Math}\label{sec_math}
One of the core strengths of \LaTeX{} is the typesetting of mathematical expressions. While the syntax takes some getting used to, you will soon discover its large economies to scale. Formulae are placed between the dollar sign operator, such as e.g. \verb+$\sum x^{2}_{i}$+ produces $\sum x^{2}_{i}$. Similarly, the mean of a discrete random variable $X$ with realisations $x_{1}, x_{2}, \ldots, x_{N}$ can be thought of as the probability-weighted sum of its realisation, i.e.

$$ \mu_{X} = \sum^{N}_{i=1}p_{i}x_{i}.$$

We can also write this same formula as an in-line text, $\mu_{X} = \sum^{N}_{i=1}p_{i}x_{i}$, without having to fret over possible implications for required spacing and formatting changes.

More complex formulae are -- with a little effort -- equally straightforward to typeset. For example, the probability density function for the normal distribution of a continuous random variable is

$$f(x) = \frac{1}{\sigma\sqrt{2\pi}}e^{-\frac{1}{2}\left(\frac{x-\mu}{\sigma}\right)^{2}}$$

You will be amazed how easy this is, once you practice a little. Instead of using the \verb+$+ or \verb+$$+ operators to invoke the math environment, we can also use the \verb+eqnarray+ environment for larger expressions. For example, we can define the mean square error (MSE) of an estimator $\hat{\theta}$ as

\begin{eqnarray}
\mathbb{E}\left[(\hat{\theta}-\theta)^2\right] &=& \mathbb{E}\left[(\hat{\theta}- \mu_{\hat{\theta}} + \mu_{\hat{\theta}} -\theta)^2\right]\\
&=&\mathbb{E}\left[(\hat{\theta}- \mu_{\hat{\theta}})^{2} + 2(\hat{\theta} - \mu_{\hat{\theta}})(\mu_{\hat{\theta}} - \theta) + (\mu_{\hat{\theta}} -\theta)^2\right]\nonumber\\
&=&\mathbb{E}\left[(\hat{\theta}- \mu_{\hat{\theta}})^{2}\right] + 2(\mu_{\hat{\theta}} - \theta)\mathbb{E}\left[\hat{\theta} - \mu_{\hat{\theta}}\right] + \mathbb{E}\Big[(\mu_{\hat{\theta}} -\theta)^2\Big]\label{eqn_mse1}\\
&=&\underbrace{\mathbb{E}\left[(\hat{\theta}- \mu_{\hat{\theta}})^{2}\right]}_{\textrm{variance of estimator } \hat{\theta}} + \underbrace{\mathbb{E}\Big[(\mu_{\hat{\theta}} -\theta)^2\Big]}_{\textrm{bias}^{2}}\label{eqn_mse2}
\end{eqnarray}

Notice that moving from equation~\eqref{eqn_mse1} to equation~\eqref{eqn_mse2}, the middle term vanishes because $\mathbb{E}\big[\hat{\theta} - \mu_{\hat{\theta}}\big]=\mathbb{E}\big[\hat{\theta}\big]-\mu_{\hat{\theta}}=0$. Thus, when choosing between two estimators, the MSE corresponds to a loss function that trades off the \emph{efficiency} of an estimator (its variance) against its \emph{unbiasedness}.

As you are stepping through the code that produces the above set of equations, you will notice another important aspect of \LaTeX{}; its internal reference and label handling. Objects such as equations, tables, figures or entire section headings can be referenced with ease by using the \verb+\label{+\emph{label\_name}\verb+}+ command first and then inserting the \verb+\ref{+\emph{label\_name}\verb+}+ or \verb+\eqref{+\emph{label\_name}\verb+}+ commands. This means that you will not have to worry about updating these references if your document changes or if you decide to apply a different numbering style (e.g. using roman numerals for section headings). The system will keep track of all labels and apply the appropriate numbering that is consistent with your formatting instructions. For example, this section is preceded by section~\ref{sec_principles} and followed by section~\ref{sec_tables}.


\subsection{Some Tables}\label{sec_tables}
The following provides a set of summary statistics for the data. This is presented in tables~\ref{tbl_Sum} and \ref{tbl_SumPlus}.

%%%%%%%%%%%%%%%%%%%%%%%%%%%%%%%%%%%%%%%%%%%%%%%%%%%%%%%%%%%%%%%%%%%%%
% Using a library such as "stargazer" or "xtable", you can
% simply use R to produce the LaTeX code for tables and then copy
% and paste that code here (NB: When overwriting the existing code
% you need to re-define the "\label{}" commands as they are not
% automatically produced by the R output).
%%%%%%%%%%%%%%%%%% TABLE CODE BEGINS %%%%%%%%%%%%%%%%%%%%%%%%%%%%%%%%
\begin{table}[htb] \centering
  \caption{Summary Statistics}
\footnotesize

\begin{tabular}{@{\extracolsep{5pt}}l c c c c c }\label{tbl_Sum}
\\[-1.8ex]\hline
\hline \\[-1.8ex]
Statistic & \multicolumn{1}{c}{\textit{N}} & \multicolumn{1}{c}{\textit{Mean}} & \multicolumn{1}{c}{\textit{St. Dev.}} & \multicolumn{1}{c}{\textit{Min}} & \multicolumn{1}{c}{\textit{Max}} \\
\hline \\[-1.8ex]
res\_vac & 1,554 & 103.627 & 113.150 & 0 & 871 \\
bus\_vac & 1,554 & 18.823 & 27.070 & 0 & 332 \\
oth\_vac & 1,554 & 0.045 & 0.297 & 0 & 7 \\
res\_tot & 1,554 & 1,525.257 & 732.713 & 0 & 7,563 \\
\hline \\[-1.8ex]
\normalsize
\end{tabular}
\end{table}
%%%%%%%%%%%%%%%%%% TABLE CODE ENDS %%%%%%%%%%%%%%%%%%%%%%%%%%%%%%%%


%%%%%%%%%%%%%%%%%% TABLE CODE BEGINS %%%%%%%%%%%%%%%%%%%%%%%%%%%%%%%%
\begin{table}[htp] \centering
  \caption{Summary Table with Additional Statistics}
\footnotesize

\begin{tabular}{@{\extracolsep{5pt}}l c c c c c c c c }\label{tbl_SumPlus}
\\[-1.8ex]\hline
\hline \\[-1.8ex]
Statistic & \multicolumn{1}{c}{\textit{N}} & \multicolumn{1}{c}{\textit{Mean}} & \multicolumn{1}{c}{\textit{St. Dev.}} & \multicolumn{1}{c}{\textit{Min}} & \multicolumn{1}{c}{\textit{Pctl(25)}} & \multicolumn{1}{c}{\textit{Median}} & \multicolumn{1}{c}{\textit{Pctl(75)}} & \multicolumn{1}{c}{\textit{Max}} \\
\hline \\[-1.8ex]
pop\_w & 1,554 & 2,501.398 & 1,787.798 & 0 & 1,123 & 2,473.5 & 3,569 & 11,632 \\
pop\_black & 1,554 & 732.153 & 1,109.979 & 0 & 33 & 173.5 & 953.5 & 5,687 \\
pop\_latino & 1,554 & 121.521 & 348.133 & 0 & 14 & 54 & 118 & 5,273 \\
hh\_median & 1,542 & 55,331.710 & 27,352.120 & 6,914 & 34,679 & 52,486.5 & 68,516.5 & 171,683 \\
\hline \\[-1.8ex]
\normalsize
\end{tabular}
\end{table}
%%%%%%%%%%%%%%%%%% TABLE CODE ENDS %%%%%%%%%%%%%%%%%%%%%%%%%%%%%%%%

\subsection{Some Figures}
Adding figures also does not take more than a couple of lines of code, using the \verb+\includegraphics{}+ command in the \verb+table+ environment. Best of all, once you have set this up in one document, you can simply adapt it for other documents. For example, figure~\ref{fig_Diversification} illustrates a positive correlation between city density and an index of relative industry diversity, calculated on the basis of all private employment at the two-digit NAICS level in 2010.

As you can see, the density-diversity relationship somewhat attenuates at the upper end of the city size distribution. This is partly because the population density in most larger cities is not exclusively driven by productivity differentials, but also by non-tradable amenity differences that both influence the quality-of-life and the quality of the business environment.

%%%%%%%%%%%%%%%%%% FIGURE CODE BEGINS %%%%%%%%%%%%%%%%%%%%%%%%%%%%%%
\begin{figure}[htp]
\caption{Relative industry diversity and population density of U.S. metropolitan areas, 2010}\label{fig_Diversification}
\begin{center}
\includegraphics[width=0.65\textwidth]{"Figures/RDI_PopDens_2010"}
\end{center}
\footnotesize
\parbox[b]{\textwidth}{\emph{Note}: The relative-diversity index (RDI) is computed by summing for each city $i$, over all $j$ sectors, the absolute value of the difference between each sector's share $s_{ij}$ in local employment and its share in national employment $s_{j}$. Formally, this yields $RDI_{i} = 1/ |\sum_{j}s_{ij} - s_{j}|$. Population density in thousands of inhabitants per square mile is plotted on a logarithmic scale.} \normalsize
\end{figure}
%%%%%%%%%%%%%%%%%% FIGURE CODE ENDS %%%%%%%%%%%%%%%%%%%%%%%%%%%%%%%%


\subsubsection{A Word on Placement}
Another nice thing about typesetting a document in this manner is that we do not have to worry about the specific placements of tables and figures. Notice that the expressions \verb+\begin{table}+ or \verb+\begin{figure}+ are followed by \verb+[htp]+, the so-called ``placement specifier'', which gives a preference for where the object, also called a ``float'', should be placed.\footnote{\texttt{[h]} means place the float where occurs in the text. \texttt{[t]} is the position at the top of the page; \texttt{[b]} denotes the bottom of the page, and \texttt{[p]} puts on a special page for floats only.} Indeed, depending on the amount of tables, figures and the general formatting constraints of your template, \LaTeX{} will figure out the optimal pagination.


\subsubsection{Some Additional Deep Insight}
Here we could have some additional text.


\section{Managing References}
In the \LaTeX{} environment, \href{http://www.bibtex.org/}{\textsc{Bib}\TeX{}} is the specific style for formatting lists of references and a \textsc{Bib}\TeX{} database is stored as a \texttt{.bib} file. It is a plain text file with a simple structure which can be viewed and edited easily. The file ``\texttt{References.bib}'' which you will find in the project folder ``\texttt{files}'' is an example of a minimal \textsc{Bib}\TeX{}.

\textsc{Bib}\TeX{} defines different entry types, with reports (\verb+@TECHREPORT{}+), articles (\verb+@ARTICLE{}+), books (\verb+@BOOK{}+) and book chapters (\verb+@INBOOK{}+) being the most common. While \textsc{Bib}\TeX{} files are simple enough to maintain using a text editor, it is more convenient to manage large \texttt{.bib} files with many entries via a reference management software that has a graphic user interface. Users of MS Word might be most familiar with the EndNotes software from Thompson Reuters for which a student version only costs \$104.99. I prefer to use software that is completely free, such as \href{http://www.mendeley.com/}{Mendeley} or, better still, \href{http://jabref.sourceforge.net/}{JabRef}, the system of choice in the open source community.


\subsection{In-text Citations}
In-text citations are really easy; we can refer to a specific piece of literature by simply inserting the respective labels for an entry that we have in our reference management database. The command \verb+\citet{+\emph{authorkey}\verb+}+ generates the most simple in-text citation. A variant of this command, \verb+\citep{}+, puts everything into parenthesis.

Notice that the specific look of the in-text citation depends on the citation style which is controlled with the command \verb+\bibliographystyle{+\emph{citationstyle}\verb+}+. The current documents uses the citation style \emph{econometrica}, the official citation style of the Econometric Society's flagship journal.\footnote{A comprehensive list of bibliography styles with examples can be found \href{http://www.mackichan.com/index.html?techtalk/632.htm~mainFrame}{here}.} The list of references is then automatically generated by the command \verb+\bibliography{+\emph{bibfile}\verb+}+ where the argument \emph{bibfile} is the name of your \texttt{.bib} file.

So, for example, I hope that you find the textbook by \citet{Dougherty2011} helpful. The article by \citet{Briant2010} might contain a surprise or two. \citet{Carrigan2012} and \citet{Goldstein2013} are useful reviews of the literature on the origins of the financial crisis.


\section{Conclusion}
No document is complete without it. Notice that there is a page break after the bibliography.

\newpage
\begin{subappendices}
\section{Additional \LaTeX{} Resources}\label{sec_resources}
There are a number of beginners' guides that I have posted on CTools. See \citet{Oetiker2011} for what is perhaps the most comprehensive introductory tutorial. Furthermore, a simple internet search with the key words ``\texttt{introduction}'' and ``\texttt{latex}'' will yield a plenty of accessible resources. Personally, I find \href{http://www.youtube.com/playlist?list=PLF975D9D3C9B50FF7}{this series} of Youtube videos particularly intuitive.

\section{My R Code}\label{sec_Rcode}
All of your homework submissions should contain the relevant \textsf{R} code that produces the key results in your work, such as tables, figures or estimation output. The \texttt{listings} package allows you to add non-formatted text as you would do with \verb+\begin{verbatim}+ but its main aim is to include the source code of any programming language within your document. The \texttt{listings} package supports highlighting of all the most common languages and it is highly customizable (see the coding set up in the preamble of the input file for this document).
\bigskip
% Copy and paste a cleaned-up version of your R code (plus relevant output) here, i.e. inside of the
% verbatim statements.
\begin{lstlisting}
setwd("[your directory structure]/UP_504/Data")

census <- read.csv("DA2_Census_TractsLarge.csv", header=T, na.strings="#N/A")
sub.census <- subset(census, census$metro>0)

library(stargazer)
library(xtable)

stargazer(sub.census[,2:7], title="Summary statistics")
stargazer(sub.census[c("pop_w", "pop_black", "pop_latino", "hh_median")],
    title="Summary table with additional statistics", median=T, iqr=T)

## Demonstration of longtable support
## Remember to insert \usepackage{longtable} on your LaTeX preamble

tbl.big <- xtable(sub.census[1:80,2:7],label="tbl.big",caption="Example of
    a Long Table Spanning Several Pages")
print(tbl.big,tabular.environment="longtable",include.colnames=T, floating=F)

## Demonstration of sidewaystable support
## Remember to insert \usepackage{rotating} on your LaTeX preamble

tbl.wide <- xtable(sub.census[1:15,2:12],label="tbl.wide",caption="A
    Wide Table")
print(tbl.wide,floating.environment="sidewaystable")
\end{lstlisting}
\normalsize

\section{A Very Long Table}
In some instances, you might need to display a large amount of data in one table that can span multiple pages or you might have a table that needs to be displayed in landscape format. The \texttt{longtable} environment is ideal for achieving the former and the \texttt{sidewaystable} environment deals with the latter task. Both functionalities are integrated nicely into the \texttt{xtable} library in \textsf{R}.

All the \textsf{R} code that was used to produce the tables in this document is printed in section~\ref{sec_Rcode} of this appendix.


%%%%%%%%%%%%%%%%%% TABLE CODE BEGINS %%%%%%%%%%%%%%%%%%%%%%%%%%%%%%%%
\begin{longtable}{rrrrrrr}
\caption{Example of a Long Table Spanning Several Pages}\label{tbl.big}\\
\hline
\hline
 & res\_vac & bus\_vac & oth\_vac & res\_tot & bus\_tot & oth\_tot \\
  \hline
\endfirsthead
 & res\_vac & bus\_vac & oth\_vac & res\_tot & bus\_tot & oth\_tot \\
  \hline
\endhead
28799 & 136 &   3 &   0 & 1307 &  36 &  27 \\
  28800 & 278 &   9 &   0 & 1216 &  44 &   2 \\
  28801 & 329 &  12 &   0 & 1480 &  38 &  14 \\
  28802 & 254 &  14 &   0 & 1144 &  56 &   2 \\
  28803 & 160 &  15 &   0 & 926 &  64 &   5 \\
  28804 & 353 &  14 &   0 & 1487 &  42 &   8 \\
  28805 & 302 &  19 &   0 & 1807 &  83 &  11 \\
  28806 & 267 &  14 &   0 & 824 &  43 &  15 \\
  28807 & 583 &  17 &   0 & 2103 &  77 &   4 \\
  28808 & 462 &  15 &   0 & 1504 &  41 &  24 \\
  28809 & 345 &  10 &   0 & 1517 &  72 &  70 \\
  28810 & 249 &   6 &   0 & 1518 &  23 &   2 \\
  28811 & 211 &  26 &   0 & 1512 & 102 &  14 \\
  28812 & 207 &  23 &   0 & 792 &  53 &  37 \\
  28813 & 526 &  11 &   0 & 1210 &  67 &  71 \\
  28814 & 273 &  17 &   0 & 2240 &  85 &  88 \\
  28815 & 218 &  14 &   0 & 1165 &  81 &  28 \\
  28816 & 259 &   7 &   0 & 1165 &  28 &   6 \\
  28817 & 145 &   7 &   0 & 1199 &  38 &   8 \\
  28818 & 219 &  18 &   0 & 896 &  88 &  12 \\
  28819 &   8 &  23 &   0 &  22 &  77 &   2 \\
  28820 & 489 &  37 &   0 & 1687 &  84 &  33 \\
  28821 & 255 &  26 &   0 & 1157 &  96 &  14 \\
  28822 & 167 &  22 &   0 & 1338 &  71 &  12 \\
  28823 &   7 &   1 &   0 & 698 &  23 &  51 \\
  28824 & 368 &  25 &   0 & 1609 &  89 &  21 \\
  28825 & 186 &  17 &   0 & 1762 &  86 &   6 \\
  28826 & 308 & 170 &   1 & 1129 & 854 & 242 \\
  28827 & 117 &  10 &   1 & 1297 &  49 &  98 \\
  28828 & 114 &  11 &   1 & 1648 &  57 &  39 \\
  28829 & 236 &  54 &   0 & 1238 & 266 & 101 \\
  28830 & 214 &  10 &   0 & 1268 &  77 &  15 \\
  28831 &  63 &  15 &   0 & 839 &  55 &  35 \\
  28832 & 373 &  21 &   0 & 1359 & 101 & 137 \\
  28833 & 159 &  20 &   0 & 1401 &  83 &  49 \\
  28834 & 277 &  23 &   0 & 2446 & 148 &  19 \\
  28835 & 267 &  26 &   0 & 1404 &  89 &  49 \\
  28836 & 157 &  25 &   0 & 771 & 122 &  30 \\
  28837 & 251 &  20 &   0 & 2511 & 113 &  46 \\
  28838 & 279 &  32 &   0 & 1745 & 102 &  14 \\
  28839 &   0 &   4 &   0 &  20 &  27 &   3 \\
  28840 &  68 &  25 &   0 & 1392 & 169 & 128 \\
  28841 &   6 &   0 &   0 & 565 &  41 &  22 \\
  28842 &  18 &   5 &   0 & 1117 &  55 &  22 \\
  28843 &  56 &  13 &   0 & 1158 & 105 &  43 \\
  28844 &  22 &   0 &   0 & 1255 &  48 &   8 \\
  28845 &  33 &   0 &   0 & 644 &  11 &   6 \\
  28846 &   1 &   0 &   0 & 767 &  29 &   7 \\
  28847 &   2 &   1 &   0 & 756 &  30 &   4 \\
  28848 & 336 &  34 &   0 & 1848 & 105 &  44 \\
  28849 & 231 &   8 &   0 & 2043 &  50 &  52 \\
  28850 &  78 &  16 &   0 & 1562 & 152 &  18 \\
  28851 &  30 &  24 &   0 & 1721 & 296 &  66 \\
  28852 &   0 &   0 &   0 & 986 &  30 &   5 \\
  28853 &  21 &   1 &   0 & 1917 &  75 &  11 \\
  28854 & 143 &  28 &   0 & 2611 & 187 &  89 \\
  28855 &  25 &   4 &   0 & 1202 &  43 &  14 \\
  28856 &  22 &   0 &   0 & 2302 &  27 &   5 \\
  28857 &  11 &   0 &   0 & 2214 &  17 &   0 \\
  28858 &  45 &   2 &   0 & 2111 &  23 &   6 \\
  28859 &   6 &   0 &   0 & 1153 &  65 &  32 \\
  28860 & 398 &  63 &   0 & 2103 & 229 &  33 \\
  28861 &  98 &  30 &   0 & 2261 & 318 & 175 \\
  28862 & 158 &  57 &   0 & 1852 & 404 & 137 \\
  28863 & 230 &  48 &   1 & 1918 & 244 & 189 \\
  28864 &  88 & 130 &   0 & 1528 & 842 & 135 \\
  28865 & 212 &  23 &   0 & 2681 & 151 &  20 \\
  28866 &  16 &  36 &   0 & 4280 & 277 &  83 \\
  28867 &   0 &   0 &   0 & 3154 & 142 & 201 \\
  28868 &  88 &  40 &   0 & 3177 & 269 & 221 \\
  28869 &  50 &  18 &   0 & 1096 & 168 &  98 \\
  28870 &  65 &   9 &   0 & 2697 & 148 & 261 \\
  28871 & 125 &  29 &   0 & 2672 & 221 &  36 \\
  28872 &  70 &   0 &   0 & 2578 &  25 &   4 \\
  28873 &   0 &   1 &   0 & 2465 &  22 &  80 \\
  28874 & 295 &  38 &   0 & 2903 & 194 &  34 \\
  28875 & 151 &  18 &   0 & 1566 & 153 &  32 \\
  28876 &  66 &  40 &   0 & 1495 & 233 &  24 \\
  28877 &  23 &   2 &   0 & 1813 &  16 &   0 \\
  28878 &  72 &  20 &   0 & 1496 & 135 &  32 \\
   \hline
\hline
\end{longtable}
%%%%%%%%%%%%%%%%%% TABLE CODE ENDS %%%%%%%%%%%%%%%%%%%%%%%%%%%%%%%%%%


%%%%%%%%%%%%%%%%%% TABLE CODE BEGINS %%%%%%%%%%%%%%%%%%%%%%%%%%%%%%%%
\begin{sidewaystable}[hbt]
\begin{center}
\begin{tabular}{rrrrrrrrrrrr}
  \hline
  \hline
 & res\_vac & bus\_vac & oth\_vac & res\_tot & bus\_tot & oth\_tot & avg\_vac\_r & avg\_vac\_b & avg\_vac\_o & metro & hh \\
  \hline
28799 & 136 &   3 &   0 & 1307 &  36 &  27 & 453.03 & 621.67 & 0.00 &   3 & 1218 \\
  28800 & 278 &   9 &   0 & 1216 &  44 &   2 & 758.30 & 823.33 & 0.00 &   3 & 898 \\
  28801 & 329 &  12 &   0 & 1480 &  38 &  14 & 493.99 & 668.83 & 0.00 &   3 & 1413 \\
  28802 & 254 &  14 &   0 & 1144 &  56 &   2 & 888.50 & 768.64 & 0.00 &   3 & 666 \\
  28803 & 160 &  15 &   0 & 926 &  64 &   5 & 769.26 & 1095.93 & 0.00 &   3 & 796 \\
  28804 & 353 &  14 &   0 & 1487 &  42 &   8 & 732.24 & 700.36 & 0.00 &   3 & 1251 \\
  28805 & 302 &  19 &   0 & 1807 &  83 &  11 & 439.34 & 810.47 & 0.00 &   3 & 1376 \\
  28806 & 267 &  14 &   0 & 824 &  43 &  15 & 782.27 & 1317.64 & 0.00 &   3 & 610 \\
  28807 & 583 &  17 &   0 & 2103 &  77 &   4 & 854.68 & 1177.06 & 0.00 &   3 & 1660 \\
  28808 & 462 &  15 &   0 & 1504 &  41 &  24 & 819.23 & 1138.00 & 0.00 &   3 & 1069 \\
  28809 & 345 &  10 &   0 & 1517 &  72 &  70 & 813.63 & 829.90 & 0.00 &   3 & 1078 \\
  28810 & 249 &   6 &   0 & 1518 &  23 &   2 & 539.71 & 1107.83 & 0.00 &   3 & 1346 \\
  28811 & 211 &  26 &   0 & 1512 & 102 &  14 & 483.38 & 1063.77 & 0.00 &   3 & 1425 \\
  28812 & 207 &  23 &   0 & 792 &  53 &  37 & 716.20 & 721.35 & 0.00 &   3 & 562 \\
  28813 & 526 &  11 &   0 & 1210 &  67 &  71 & 967.73 & 907.46 & 0.00 &   3 & 778 \\
   \hline
\end{tabular}
\caption{A Wide Table}
\label{tbl.wide}
\end{center}
\end{sidewaystable}
%%%%%%%%%%%%%%%%%% TABLE CODE ENDS %%%%%%%%%%%%%%%%%%%%%%%%%%%%%%%%%%

\section{Step-by-step installation instructions}
\begin{enumerate}
\item Visit \href{https://www.sharelatex.com}{\texttt{https://www.sharelatex.com}}, create an account and log in.
\item On CTools, go to the folder \texttt{UP 504 002 W13} $\rightarrow$ \texttt{Resources} $\rightarrow$ \texttt{LaTeX} $\rightarrow$ \texttt{Programs} and download the file \texttt{UP504\_LaTeXProject.zip} to a location on your machine that you can find again. Do not unzip the file.
\item In your \textsc{Share}\LaTeX{} environment, select the ``Projects'' item in the menu and click on the green \emph{New Project} button. Select ``Upload Zipped Project'' from the drop list and select the \texttt{.zip} folder that you have downloaded from CTools. Upload that folder.
\item Once the folder has been uploaded, you will see a screen with a section ``Collaborators'' and a section ``Options''. In the ``Root Document'' field in the ``Options'' section, select the \texttt{UP504\_HomeworkTemplate.tex}.
\item Click on the \texttt{UP504\_HomeworkTemplate.tex} icon in the vertical navigation pane on the left. You should now see the context of the input file. This is the file that gives \LaTeX{} all the necessary formatting instructions.
\item You can now compile this file into a \texttt{.pdf} output file by clicking on the green \emph{Recompile} button. Your compiled document is now visible. This file serves both as a template for future homeworks and as a manual that introduces the most essential functions for our class.
\item For future homeworks, you either simply replace the text in the template file or you can create a new blank project, copying and pasting the necessary code from the template.
\end{enumerate}

\renewcommand\bibname{References}
\bibliographystyle{econometrica}
\bibliography{Literature/OverLord}
\IfFileExists{Literature/IntroLaTeX.bbl}{\chapter{A Concise Introduction to \LaTeX}\label{chap_LaTeX}\index{software!\LaTeX}

\section{Getting Started}
This template is intended to facilitate the adoption of \LaTeX{} for future homework submissions in UP504. This \texttt{.tex} file has been set up such that you can begin producing fairly advanced documents without getting too bogged down with coding (see appendix~\ref{sec_resources} for additional \LaTeX{} resources).

\subsection{Some Motivation}\label{sec_motivation}
The many advantages of using a typesetting system like \LaTeX{} arise from the fact that the computer takes care of all formatting issues, freeing you up to concentrate on ``the highest and best use'' of your time: delivering meaningful content. Unlike a conventional WYSIWYG\footnote{In WYSIWYG (What you see is what you get) text editing systems, the input and final output file are identical and are simultaneously displayed on screen. As you can see, footnotes are really easy too!} word processor such as MS Word, \LaTeX{} does not distract you by a clever reverse delegation of sectioning, pagination, object placement, numbering and other formatting tasks. Imagine being able to convert a letter-sized, one-column document with footnotes, tables and figures into an A4-sized, two-column layout with double-spaced, numbered lines and endnotes in less than 30 seconds! With \LaTeX{}, all you need is three easy commands and a few parameter changes to achieve that. The commands are \verb+\let\footnote=\endnote+, \verb+\doublespace+, and \verb+\linenumbers+ plus adding the terms \verb+twocolumn+ and \verb+a4+ to the \verb+\documentclass{}+ command.

When it comes to producing a professionally typeset document, there are of course many additional advantages to using \LaTeX{}, such as kerning (the proportional spacing between characters), the use of common ligatures (joining of one or more letters into a single glyph) or advanced hyphenation. You can find a comprehensive illustration of these aspects \href{http://nitens.org/taraborelli/latex}{here}.

\subsection{Some Basics}\label{sec_basics}
The input for \LaTeX{} is a plain text file that contains the text of the document, as well as the commands that tell \LaTeX{}
how to typeset the text. ``Whitespace'' characters, such as blanks or tabs, are treated uniformly as ``space'', several consecutive whitespace characters are treated as one space. \LaTeX{} treats a single line break as whitespace and an empty line between two lines of text defines the end of a paragraph. Several empty lines are treated the same as one empty line.

\subsubsection{Command Syntax}
\LaTeX{} commands are case sensitive, and take one of the following two formats:

\begin{itemize}
\item They start with a backslash \verb+\+ and then have a name consisting of letters only. Command names are terminated by a space, a number or any other `non-letter' character.
\item Alternatively, commands consist of a backslash and exactly one non-letter character.
\end{itemize}

In many instances, both command formats are available to achieve the same thing. For example, \verb+\newline+ or \verb+\\+ starts a new line without starting a new paragraph. \verb+\\*+ additionally prohibits a page break after the forced line break. Some commands require a parameter which has to be given between curly brackets \verb+{...}+ after the command name. Some commands take optional parameters which are inserted after the command name in square brackets \verb+[...]+. Generally, this implies the following command syntax:\\

\begin{center}
\verb+\+\texttt{command}\verb+[+\emph{optional parameter}\verb+]{+\emph{parameter}\verb+}+.
\end{center}


\subsubsection{Document Structure}
In order for \LaTeX{} input files to be compiled successfully, they have to follow a certain structure. Thus every input file must start with the command \verb+\documentclass{...}+ which defines what sort of document you intend to write. After that, you can define
commands to influence the style of the whole document, or load packages that add new features to the \LaTeX{} system. To load such a package you use the command \verb+\usepackage{...}+.

When all the setup work is done, you start the body of the text with the command \verb+\begin{document}+. Now you enter the text mixed with some useful \LaTeX{} commands. The area between \verb+\documentclass+ and \verb+\begin{document}+ is called the \emph{preamble}. The end of every input file is indicated by the \verb+\end{document}+ command. Anything that follows this command will be ignored by the \LaTeX{} compiler. For example, the simplest possible \LaTeX{} document is generated by the following code:

\begin{verbatim}
\documentclass{article}
\begin{document}
Small is beautiful.
\end{document}
\end{verbatim}

More info on document structure can be found \href{http://en.wikibooks.org/wiki/LaTeX/Document_Structure}{here}.

\subsection{Some Principles}\label{sec_principles}
When it comes to coding, \LaTeX{} is no different from any other interpreted language. Successful coding has a high fixed cost, a steep learning curve, but yields large economies to scale. Most questions about code are made more answerable by the addition of a minimal working example (MWE) or a short, self-contained, correct example (SSCCE). Your debugging question may have an obvious answer, or it may not. When enlisting the help of others in debugging your code, think about it from the point of view of the person whose help you're seeking. Thus you should always provide a block of code that can be copied, pasted directly into a file, compiled, tweaked, corrected, and pasted back.

In this context, ``minimal'' means that the problem is isolated and there is nothing in the sample that distracts from the error you have or the effect you desire. ``Working'' doesn't mean the problem is solved, it means the code sample can be copied-and-pasted and directly compiled.\footnote{The site \href{http://www.sscce.org}{sscce.org} is dedicated to the topic. There are also several articles on the topic as it applies specifically to LaTeX: \href{http://www.tex.ac.uk/cgi-bin/texfaq2html?label=minxampl}{``How to make a `minimum example'"} from the UK \TeX{} FAQ, \href{http://www.minimalbeispiel.de/mini-en.html}{``What is a minimal working example?"} from the \texttt{de.comp.text.tex} newsgroup.}

\subsection{Some Math}\label{sec_math}
One of the core strengths of \LaTeX{} is the typesetting of mathematical expressions. While the syntax takes some getting used to, you will soon discover its large economies to scale. Formulae are placed between the dollar sign operator, such as e.g. \verb+$\sum x^{2}_{i}$+ produces $\sum x^{2}_{i}$. Similarly, the mean of a discrete random variable $X$ with realisations $x_{1}, x_{2}, \ldots, x_{N}$ can be thought of as the probability-weighted sum of its realisation, i.e.

$$ \mu_{X} = \sum^{N}_{i=1}p_{i}x_{i}.$$

We can also write this same formula as an in-line text, $\mu_{X} = \sum^{N}_{i=1}p_{i}x_{i}$, without having to fret over possible implications for required spacing and formatting changes.

More complex formulae are -- with a little effort -- equally straightforward to typeset. For example, the probability density function for the normal distribution of a continuous random variable is

$$f(x) = \frac{1}{\sigma\sqrt{2\pi}}e^{-\frac{1}{2}\left(\frac{x-\mu}{\sigma}\right)^{2}}$$

You will be amazed how easy this is, once you practice a little. Instead of using the \verb+$+ or \verb+$$+ operators to invoke the math environment, we can also use the \verb+eqnarray+ environment for larger expressions. For example, we can define the mean square error (MSE) of an estimator $\hat{\theta}$ as

\begin{eqnarray}
\mathbb{E}\left[(\hat{\theta}-\theta)^2\right] &=& \mathbb{E}\left[(\hat{\theta}- \mu_{\hat{\theta}} + \mu_{\hat{\theta}} -\theta)^2\right]\\
&=&\mathbb{E}\left[(\hat{\theta}- \mu_{\hat{\theta}})^{2} + 2(\hat{\theta} - \mu_{\hat{\theta}})(\mu_{\hat{\theta}} - \theta) + (\mu_{\hat{\theta}} -\theta)^2\right]\nonumber\\
&=&\mathbb{E}\left[(\hat{\theta}- \mu_{\hat{\theta}})^{2}\right] + 2(\mu_{\hat{\theta}} - \theta)\mathbb{E}\left[\hat{\theta} - \mu_{\hat{\theta}}\right] + \mathbb{E}\Big[(\mu_{\hat{\theta}} -\theta)^2\Big]\label{eqn_mse1}\\
&=&\underbrace{\mathbb{E}\left[(\hat{\theta}- \mu_{\hat{\theta}})^{2}\right]}_{\textrm{variance of estimator } \hat{\theta}} + \underbrace{\mathbb{E}\Big[(\mu_{\hat{\theta}} -\theta)^2\Big]}_{\textrm{bias}^{2}}\label{eqn_mse2}
\end{eqnarray}

Notice that moving from equation~\eqref{eqn_mse1} to equation~\eqref{eqn_mse2}, the middle term vanishes because $\mathbb{E}\big[\hat{\theta} - \mu_{\hat{\theta}}\big]=\mathbb{E}\big[\hat{\theta}\big]-\mu_{\hat{\theta}}=0$. Thus, when choosing between two estimators, the MSE corresponds to a loss function that trades off the \emph{efficiency} of an estimator (its variance) against its \emph{unbiasedness}.

As you are stepping through the code that produces the above set of equations, you will notice another important aspect of \LaTeX{}; its internal reference and label handling. Objects such as equations, tables, figures or entire section headings can be referenced with ease by using the \verb+\label{+\emph{label\_name}\verb+}+ command first and then inserting the \verb+\ref{+\emph{label\_name}\verb+}+ or \verb+\eqref{+\emph{label\_name}\verb+}+ commands. This means that you will not have to worry about updating these references if your document changes or if you decide to apply a different numbering style (e.g. using roman numerals for section headings). The system will keep track of all labels and apply the appropriate numbering that is consistent with your formatting instructions. For example, this section is preceded by section~\ref{sec_principles} and followed by section~\ref{sec_tables}.


\subsection{Some Tables}\label{sec_tables}
The following provides a set of summary statistics for the data. This is presented in tables~\ref{tbl_Sum} and \ref{tbl_SumPlus}.

%%%%%%%%%%%%%%%%%%%%%%%%%%%%%%%%%%%%%%%%%%%%%%%%%%%%%%%%%%%%%%%%%%%%%
% Using a library such as "stargazer" or "xtable", you can
% simply use R to produce the LaTeX code for tables and then copy
% and paste that code here (NB: When overwriting the existing code
% you need to re-define the "\label{}" commands as they are not
% automatically produced by the R output).
%%%%%%%%%%%%%%%%%% TABLE CODE BEGINS %%%%%%%%%%%%%%%%%%%%%%%%%%%%%%%%
\begin{table}[htb] \centering
  \caption{Summary Statistics}
\footnotesize

\begin{tabular}{@{\extracolsep{5pt}}l c c c c c }\label{tbl_Sum}
\\[-1.8ex]\hline
\hline \\[-1.8ex]
Statistic & \multicolumn{1}{c}{\textit{N}} & \multicolumn{1}{c}{\textit{Mean}} & \multicolumn{1}{c}{\textit{St. Dev.}} & \multicolumn{1}{c}{\textit{Min}} & \multicolumn{1}{c}{\textit{Max}} \\
\hline \\[-1.8ex]
res\_vac & 1,554 & 103.627 & 113.150 & 0 & 871 \\
bus\_vac & 1,554 & 18.823 & 27.070 & 0 & 332 \\
oth\_vac & 1,554 & 0.045 & 0.297 & 0 & 7 \\
res\_tot & 1,554 & 1,525.257 & 732.713 & 0 & 7,563 \\
\hline \\[-1.8ex]
\normalsize
\end{tabular}
\end{table}
%%%%%%%%%%%%%%%%%% TABLE CODE ENDS %%%%%%%%%%%%%%%%%%%%%%%%%%%%%%%%


%%%%%%%%%%%%%%%%%% TABLE CODE BEGINS %%%%%%%%%%%%%%%%%%%%%%%%%%%%%%%%
\begin{table}[htp] \centering
  \caption{Summary Table with Additional Statistics}
\footnotesize

\begin{tabular}{@{\extracolsep{5pt}}l c c c c c c c c }\label{tbl_SumPlus}
\\[-1.8ex]\hline
\hline \\[-1.8ex]
Statistic & \multicolumn{1}{c}{\textit{N}} & \multicolumn{1}{c}{\textit{Mean}} & \multicolumn{1}{c}{\textit{St. Dev.}} & \multicolumn{1}{c}{\textit{Min}} & \multicolumn{1}{c}{\textit{Pctl(25)}} & \multicolumn{1}{c}{\textit{Median}} & \multicolumn{1}{c}{\textit{Pctl(75)}} & \multicolumn{1}{c}{\textit{Max}} \\
\hline \\[-1.8ex]
pop\_w & 1,554 & 2,501.398 & 1,787.798 & 0 & 1,123 & 2,473.5 & 3,569 & 11,632 \\
pop\_black & 1,554 & 732.153 & 1,109.979 & 0 & 33 & 173.5 & 953.5 & 5,687 \\
pop\_latino & 1,554 & 121.521 & 348.133 & 0 & 14 & 54 & 118 & 5,273 \\
hh\_median & 1,542 & 55,331.710 & 27,352.120 & 6,914 & 34,679 & 52,486.5 & 68,516.5 & 171,683 \\
\hline \\[-1.8ex]
\normalsize
\end{tabular}
\end{table}
%%%%%%%%%%%%%%%%%% TABLE CODE ENDS %%%%%%%%%%%%%%%%%%%%%%%%%%%%%%%%

\subsection{Some Figures}
Adding figures also does not take more than a couple of lines of code, using the \verb+\includegraphics{}+ command in the \verb+table+ environment. Best of all, once you have set this up in one document, you can simply adapt it for other documents. For example, figure~\ref{fig_Diversification} illustrates a positive correlation between city density and an index of relative industry diversity, calculated on the basis of all private employment at the two-digit NAICS level in 2010.

As you can see, the density-diversity relationship somewhat attenuates at the upper end of the city size distribution. This is partly because the population density in most larger cities is not exclusively driven by productivity differentials, but also by non-tradable amenity differences that both influence the quality-of-life and the quality of the business environment.

%%%%%%%%%%%%%%%%%% FIGURE CODE BEGINS %%%%%%%%%%%%%%%%%%%%%%%%%%%%%%
\begin{figure}[htp]
\caption{Relative industry diversity and population density of U.S. metropolitan areas, 2010}\label{fig_Diversification}
\begin{center}
\includegraphics[width=0.65\textwidth]{"Figures/RDI_PopDens_2010"}
\end{center}
\footnotesize
\parbox[b]{\textwidth}{\emph{Note}: The relative-diversity index (RDI) is computed by summing for each city $i$, over all $j$ sectors, the absolute value of the difference between each sector's share $s_{ij}$ in local employment and its share in national employment $s_{j}$. Formally, this yields $RDI_{i} = 1/ |\sum_{j}s_{ij} - s_{j}|$. Population density in thousands of inhabitants per square mile is plotted on a logarithmic scale.} \normalsize
\end{figure}
%%%%%%%%%%%%%%%%%% FIGURE CODE ENDS %%%%%%%%%%%%%%%%%%%%%%%%%%%%%%%%


\subsubsection{A Word on Placement}
Another nice thing about typesetting a document in this manner is that we do not have to worry about the specific placements of tables and figures. Notice that the expressions \verb+\begin{table}+ or \verb+\begin{figure}+ are followed by \verb+[htp]+, the so-called ``placement specifier'', which gives a preference for where the object, also called a ``float'', should be placed.\footnote{\texttt{[h]} means place the float where occurs in the text. \texttt{[t]} is the position at the top of the page; \texttt{[b]} denotes the bottom of the page, and \texttt{[p]} puts on a special page for floats only.} Indeed, depending on the amount of tables, figures and the general formatting constraints of your template, \LaTeX{} will figure out the optimal pagination.


\subsubsection{Some Additional Deep Insight}
Here we could have some additional text.


\section{Managing References}
In the \LaTeX{} environment, \href{http://www.bibtex.org/}{\textsc{Bib}\TeX{}} is the specific style for formatting lists of references and a \textsc{Bib}\TeX{} database is stored as a \texttt{.bib} file. It is a plain text file with a simple structure which can be viewed and edited easily. The file ``\texttt{References.bib}'' which you will find in the project folder ``\texttt{files}'' is an example of a minimal \textsc{Bib}\TeX{}.

\textsc{Bib}\TeX{} defines different entry types, with reports (\verb+@TECHREPORT{}+), articles (\verb+@ARTICLE{}+), books (\verb+@BOOK{}+) and book chapters (\verb+@INBOOK{}+) being the most common. While \textsc{Bib}\TeX{} files are simple enough to maintain using a text editor, it is more convenient to manage large \texttt{.bib} files with many entries via a reference management software that has a graphic user interface. Users of MS Word might be most familiar with the EndNotes software from Thompson Reuters for which a student version only costs \$104.99. I prefer to use software that is completely free, such as \href{http://www.mendeley.com/}{Mendeley} or, better still, \href{http://jabref.sourceforge.net/}{JabRef}, the system of choice in the open source community.


\subsection{In-text Citations}
In-text citations are really easy; we can refer to a specific piece of literature by simply inserting the respective labels for an entry that we have in our reference management database. The command \verb+\citet{+\emph{authorkey}\verb+}+ generates the most simple in-text citation. A variant of this command, \verb+\citep{}+, puts everything into parenthesis.

Notice that the specific look of the in-text citation depends on the citation style which is controlled with the command \verb+\bibliographystyle{+\emph{citationstyle}\verb+}+. The current documents uses the citation style \emph{econometrica}, the official citation style of the Econometric Society's flagship journal.\footnote{A comprehensive list of bibliography styles with examples can be found \href{http://www.mackichan.com/index.html?techtalk/632.htm~mainFrame}{here}.} The list of references is then automatically generated by the command \verb+\bibliography{+\emph{bibfile}\verb+}+ where the argument \emph{bibfile} is the name of your \texttt{.bib} file.

So, for example, I hope that you find the textbook by \citet{Dougherty2011} helpful. The article by \citet{Briant2010} might contain a surprise or two. \citet{Carrigan2012} and \citet{Goldstein2013} are useful reviews of the literature on the origins of the financial crisis.


\section{Conclusion}
No document is complete without it. Notice that there is a page break after the bibliography.

\newpage
\begin{subappendices}
\section{Additional \LaTeX{} Resources}\label{sec_resources}
There are a number of beginners' guides that I have posted on CTools. See \citet{Oetiker2011} for what is perhaps the most comprehensive introductory tutorial. Furthermore, a simple internet search with the key words ``\texttt{introduction}'' and ``\texttt{latex}'' will yield a plenty of accessible resources. Personally, I find \href{http://www.youtube.com/playlist?list=PLF975D9D3C9B50FF7}{this series} of Youtube videos particularly intuitive.

\section{My R Code}\label{sec_Rcode}
All of your homework submissions should contain the relevant \textsf{R} code that produces the key results in your work, such as tables, figures or estimation output. The \texttt{listings} package allows you to add non-formatted text as you would do with \verb+\begin{verbatim}+ but its main aim is to include the source code of any programming language within your document. The \texttt{listings} package supports highlighting of all the most common languages and it is highly customizable (see the coding set up in the preamble of the input file for this document).
\bigskip
% Copy and paste a cleaned-up version of your R code (plus relevant output) here, i.e. inside of the
% verbatim statements.
\begin{lstlisting}
setwd("[your directory structure]/UP_504/Data")

census <- read.csv("DA2_Census_TractsLarge.csv", header=T, na.strings="#N/A")
sub.census <- subset(census, census$metro>0)

library(stargazer)
library(xtable)

stargazer(sub.census[,2:7], title="Summary statistics")
stargazer(sub.census[c("pop_w", "pop_black", "pop_latino", "hh_median")],
    title="Summary table with additional statistics", median=T, iqr=T)

## Demonstration of longtable support
## Remember to insert \usepackage{longtable} on your LaTeX preamble

tbl.big <- xtable(sub.census[1:80,2:7],label="tbl.big",caption="Example of
    a Long Table Spanning Several Pages")
print(tbl.big,tabular.environment="longtable",include.colnames=T, floating=F)

## Demonstration of sidewaystable support
## Remember to insert \usepackage{rotating} on your LaTeX preamble

tbl.wide <- xtable(sub.census[1:15,2:12],label="tbl.wide",caption="A
    Wide Table")
print(tbl.wide,floating.environment="sidewaystable")
\end{lstlisting}
\normalsize

\section{A Very Long Table}
In some instances, you might need to display a large amount of data in one table that can span multiple pages or you might have a table that needs to be displayed in landscape format. The \texttt{longtable} environment is ideal for achieving the former and the \texttt{sidewaystable} environment deals with the latter task. Both functionalities are integrated nicely into the \texttt{xtable} library in \textsf{R}.

All the \textsf{R} code that was used to produce the tables in this document is printed in section~\ref{sec_Rcode} of this appendix.


%%%%%%%%%%%%%%%%%% TABLE CODE BEGINS %%%%%%%%%%%%%%%%%%%%%%%%%%%%%%%%
\begin{longtable}{rrrrrrr}
\caption{Example of a Long Table Spanning Several Pages}\label{tbl.big}\\
\hline
\hline
 & res\_vac & bus\_vac & oth\_vac & res\_tot & bus\_tot & oth\_tot \\
  \hline
\endfirsthead
 & res\_vac & bus\_vac & oth\_vac & res\_tot & bus\_tot & oth\_tot \\
  \hline
\endhead
28799 & 136 &   3 &   0 & 1307 &  36 &  27 \\
  28800 & 278 &   9 &   0 & 1216 &  44 &   2 \\
  28801 & 329 &  12 &   0 & 1480 &  38 &  14 \\
  28802 & 254 &  14 &   0 & 1144 &  56 &   2 \\
  28803 & 160 &  15 &   0 & 926 &  64 &   5 \\
  28804 & 353 &  14 &   0 & 1487 &  42 &   8 \\
  28805 & 302 &  19 &   0 & 1807 &  83 &  11 \\
  28806 & 267 &  14 &   0 & 824 &  43 &  15 \\
  28807 & 583 &  17 &   0 & 2103 &  77 &   4 \\
  28808 & 462 &  15 &   0 & 1504 &  41 &  24 \\
  28809 & 345 &  10 &   0 & 1517 &  72 &  70 \\
  28810 & 249 &   6 &   0 & 1518 &  23 &   2 \\
  28811 & 211 &  26 &   0 & 1512 & 102 &  14 \\
  28812 & 207 &  23 &   0 & 792 &  53 &  37 \\
  28813 & 526 &  11 &   0 & 1210 &  67 &  71 \\
  28814 & 273 &  17 &   0 & 2240 &  85 &  88 \\
  28815 & 218 &  14 &   0 & 1165 &  81 &  28 \\
  28816 & 259 &   7 &   0 & 1165 &  28 &   6 \\
  28817 & 145 &   7 &   0 & 1199 &  38 &   8 \\
  28818 & 219 &  18 &   0 & 896 &  88 &  12 \\
  28819 &   8 &  23 &   0 &  22 &  77 &   2 \\
  28820 & 489 &  37 &   0 & 1687 &  84 &  33 \\
  28821 & 255 &  26 &   0 & 1157 &  96 &  14 \\
  28822 & 167 &  22 &   0 & 1338 &  71 &  12 \\
  28823 &   7 &   1 &   0 & 698 &  23 &  51 \\
  28824 & 368 &  25 &   0 & 1609 &  89 &  21 \\
  28825 & 186 &  17 &   0 & 1762 &  86 &   6 \\
  28826 & 308 & 170 &   1 & 1129 & 854 & 242 \\
  28827 & 117 &  10 &   1 & 1297 &  49 &  98 \\
  28828 & 114 &  11 &   1 & 1648 &  57 &  39 \\
  28829 & 236 &  54 &   0 & 1238 & 266 & 101 \\
  28830 & 214 &  10 &   0 & 1268 &  77 &  15 \\
  28831 &  63 &  15 &   0 & 839 &  55 &  35 \\
  28832 & 373 &  21 &   0 & 1359 & 101 & 137 \\
  28833 & 159 &  20 &   0 & 1401 &  83 &  49 \\
  28834 & 277 &  23 &   0 & 2446 & 148 &  19 \\
  28835 & 267 &  26 &   0 & 1404 &  89 &  49 \\
  28836 & 157 &  25 &   0 & 771 & 122 &  30 \\
  28837 & 251 &  20 &   0 & 2511 & 113 &  46 \\
  28838 & 279 &  32 &   0 & 1745 & 102 &  14 \\
  28839 &   0 &   4 &   0 &  20 &  27 &   3 \\
  28840 &  68 &  25 &   0 & 1392 & 169 & 128 \\
  28841 &   6 &   0 &   0 & 565 &  41 &  22 \\
  28842 &  18 &   5 &   0 & 1117 &  55 &  22 \\
  28843 &  56 &  13 &   0 & 1158 & 105 &  43 \\
  28844 &  22 &   0 &   0 & 1255 &  48 &   8 \\
  28845 &  33 &   0 &   0 & 644 &  11 &   6 \\
  28846 &   1 &   0 &   0 & 767 &  29 &   7 \\
  28847 &   2 &   1 &   0 & 756 &  30 &   4 \\
  28848 & 336 &  34 &   0 & 1848 & 105 &  44 \\
  28849 & 231 &   8 &   0 & 2043 &  50 &  52 \\
  28850 &  78 &  16 &   0 & 1562 & 152 &  18 \\
  28851 &  30 &  24 &   0 & 1721 & 296 &  66 \\
  28852 &   0 &   0 &   0 & 986 &  30 &   5 \\
  28853 &  21 &   1 &   0 & 1917 &  75 &  11 \\
  28854 & 143 &  28 &   0 & 2611 & 187 &  89 \\
  28855 &  25 &   4 &   0 & 1202 &  43 &  14 \\
  28856 &  22 &   0 &   0 & 2302 &  27 &   5 \\
  28857 &  11 &   0 &   0 & 2214 &  17 &   0 \\
  28858 &  45 &   2 &   0 & 2111 &  23 &   6 \\
  28859 &   6 &   0 &   0 & 1153 &  65 &  32 \\
  28860 & 398 &  63 &   0 & 2103 & 229 &  33 \\
  28861 &  98 &  30 &   0 & 2261 & 318 & 175 \\
  28862 & 158 &  57 &   0 & 1852 & 404 & 137 \\
  28863 & 230 &  48 &   1 & 1918 & 244 & 189 \\
  28864 &  88 & 130 &   0 & 1528 & 842 & 135 \\
  28865 & 212 &  23 &   0 & 2681 & 151 &  20 \\
  28866 &  16 &  36 &   0 & 4280 & 277 &  83 \\
  28867 &   0 &   0 &   0 & 3154 & 142 & 201 \\
  28868 &  88 &  40 &   0 & 3177 & 269 & 221 \\
  28869 &  50 &  18 &   0 & 1096 & 168 &  98 \\
  28870 &  65 &   9 &   0 & 2697 & 148 & 261 \\
  28871 & 125 &  29 &   0 & 2672 & 221 &  36 \\
  28872 &  70 &   0 &   0 & 2578 &  25 &   4 \\
  28873 &   0 &   1 &   0 & 2465 &  22 &  80 \\
  28874 & 295 &  38 &   0 & 2903 & 194 &  34 \\
  28875 & 151 &  18 &   0 & 1566 & 153 &  32 \\
  28876 &  66 &  40 &   0 & 1495 & 233 &  24 \\
  28877 &  23 &   2 &   0 & 1813 &  16 &   0 \\
  28878 &  72 &  20 &   0 & 1496 & 135 &  32 \\
   \hline
\hline
\end{longtable}
%%%%%%%%%%%%%%%%%% TABLE CODE ENDS %%%%%%%%%%%%%%%%%%%%%%%%%%%%%%%%%%


%%%%%%%%%%%%%%%%%% TABLE CODE BEGINS %%%%%%%%%%%%%%%%%%%%%%%%%%%%%%%%
\begin{sidewaystable}[hbt]
\begin{center}
\begin{tabular}{rrrrrrrrrrrr}
  \hline
  \hline
 & res\_vac & bus\_vac & oth\_vac & res\_tot & bus\_tot & oth\_tot & avg\_vac\_r & avg\_vac\_b & avg\_vac\_o & metro & hh \\
  \hline
28799 & 136 &   3 &   0 & 1307 &  36 &  27 & 453.03 & 621.67 & 0.00 &   3 & 1218 \\
  28800 & 278 &   9 &   0 & 1216 &  44 &   2 & 758.30 & 823.33 & 0.00 &   3 & 898 \\
  28801 & 329 &  12 &   0 & 1480 &  38 &  14 & 493.99 & 668.83 & 0.00 &   3 & 1413 \\
  28802 & 254 &  14 &   0 & 1144 &  56 &   2 & 888.50 & 768.64 & 0.00 &   3 & 666 \\
  28803 & 160 &  15 &   0 & 926 &  64 &   5 & 769.26 & 1095.93 & 0.00 &   3 & 796 \\
  28804 & 353 &  14 &   0 & 1487 &  42 &   8 & 732.24 & 700.36 & 0.00 &   3 & 1251 \\
  28805 & 302 &  19 &   0 & 1807 &  83 &  11 & 439.34 & 810.47 & 0.00 &   3 & 1376 \\
  28806 & 267 &  14 &   0 & 824 &  43 &  15 & 782.27 & 1317.64 & 0.00 &   3 & 610 \\
  28807 & 583 &  17 &   0 & 2103 &  77 &   4 & 854.68 & 1177.06 & 0.00 &   3 & 1660 \\
  28808 & 462 &  15 &   0 & 1504 &  41 &  24 & 819.23 & 1138.00 & 0.00 &   3 & 1069 \\
  28809 & 345 &  10 &   0 & 1517 &  72 &  70 & 813.63 & 829.90 & 0.00 &   3 & 1078 \\
  28810 & 249 &   6 &   0 & 1518 &  23 &   2 & 539.71 & 1107.83 & 0.00 &   3 & 1346 \\
  28811 & 211 &  26 &   0 & 1512 & 102 &  14 & 483.38 & 1063.77 & 0.00 &   3 & 1425 \\
  28812 & 207 &  23 &   0 & 792 &  53 &  37 & 716.20 & 721.35 & 0.00 &   3 & 562 \\
  28813 & 526 &  11 &   0 & 1210 &  67 &  71 & 967.73 & 907.46 & 0.00 &   3 & 778 \\
   \hline
\end{tabular}
\caption{A Wide Table}
\label{tbl.wide}
\end{center}
\end{sidewaystable}
%%%%%%%%%%%%%%%%%% TABLE CODE ENDS %%%%%%%%%%%%%%%%%%%%%%%%%%%%%%%%%%

\section{Step-by-step installation instructions}
\begin{enumerate}
\item Visit \href{https://www.sharelatex.com}{\texttt{https://www.sharelatex.com}}, create an account and log in.
\item On CTools, go to the folder \texttt{UP 504 002 W13} $\rightarrow$ \texttt{Resources} $\rightarrow$ \texttt{LaTeX} $\rightarrow$ \texttt{Programs} and download the file \texttt{UP504\_LaTeXProject.zip} to a location on your machine that you can find again. Do not unzip the file.
\item In your \textsc{Share}\LaTeX{} environment, select the ``Projects'' item in the menu and click on the green \emph{New Project} button. Select ``Upload Zipped Project'' from the drop list and select the \texttt{.zip} folder that you have downloaded from CTools. Upload that folder.
\item Once the folder has been uploaded, you will see a screen with a section ``Collaborators'' and a section ``Options''. In the ``Root Document'' field in the ``Options'' section, select the \texttt{UP504\_HomeworkTemplate.tex}.
\item Click on the \texttt{UP504\_HomeworkTemplate.tex} icon in the vertical navigation pane on the left. You should now see the context of the input file. This is the file that gives \LaTeX{} all the necessary formatting instructions.
\item You can now compile this file into a \texttt{.pdf} output file by clicking on the green \emph{Recompile} button. Your compiled document is now visible. This file serves both as a template for future homeworks and as a manual that introduces the most essential functions for our class.
\item For future homeworks, you either simply replace the text in the template file or you can create a new blank project, copying and pasting the necessary code from the template.
\end{enumerate}

\renewcommand\bibname{References}
\bibliographystyle{econometrica}
\bibliography{Literature/OverLord}
\IfFileExists{Literature/IntroLaTeX.bbl}{\input{Literature/IntroLaTeX.bbl}}{\typeout{Need to run bibtex}}
\end{subappendices} }{\typeout{Need to run bibtex}}
\end{subappendices} }{\typeout{Need to run bibtex}}
\end{subappendices} }{\typeout{Need to run bibtex}}
\end{subappendices} 